% Lab 2 of Intelligent Robotics (ASY2): a polygon map of a room from one laser scan, by least squares and RANSAC.
% Build: make  (latexmk -pdf). Figures are TikZ files in figures/, no external graphics.
\documentclass[11pt,a4paper,parskip=half]{scrartcl}
\usepackage[english]{babel}
\usepackage[T1]{fontenc}
\usepackage[utf8]{inputenc}
\usepackage[sfdefault,lining]{FiraSans}
\usepackage[scale=0.88]{FiraMono}
\usepackage[a4paper,left=2.4cm,right=2.4cm,top=2.6cm,bottom=2.6cm,headsep=0.8cm]{geometry}
\usepackage{amsmath,amssymb}
\usepackage{booktabs,array,tabularx}
\usepackage{xcolor}
\usepackage{listings}
\usepackage{tikz}
\usetikzlibrary{arrows.meta,calc,shapes.geometric}
\usepackage[skins,breakable]{tcolorbox}
\usepackage{fontawesome5}
\usepackage[font=small,labelfont={bf,color=ohmred},skip=4pt]{caption}
\usepackage{enumitem}
\usepackage{microtype}
\usepackage[headsepline=0.4pt]{scrlayer-scrpage}
\usepackage[colorlinks=true,linkcolor=ohmred,urlcolor=ohmred]{hyperref}
\usepackage{cleveref}

% ---------------------------------------------------------------- colours and headings
\definecolor{ohmred}{HTML}{C8372D}
\definecolor{ohmgreen}{HTML}{2E8B57}
\definecolor{ohmamber}{HTML}{D98E04}
\definecolor{ohmgrey}{HTML}{5A5A5A}
\addtokomafont{disposition}{\color{ohmred}}
\addtokomafont{section}{\Large}
\addtokomafont{subsection}{\large}
\setkomafont{pageheadfoot}{\small\color{ohmgrey}}
\setkomafont{pagenumber}{\small\bfseries\color{ohmred}}
\clearpairofpagestyles
\ihead{Intelligent Robotics (ASY2) · Lab 2}
\ohead{Polygon map}
\ofoot{\pagemark}
\pagestyle{scrheadings}

\lstset{basicstyle=\ttfamily\small, columns=fullflexible, keepspaces=true,
        showstringspaces=false, breaklines=true, frame=none, backgroundcolor=\color{black!5},
        xleftmargin=8pt, xrightmargin=8pt, framexleftmargin=6pt, framexrightmargin=6pt,
        aboveskip=4pt, belowskip=4pt}

\newcommand{\thema}[1]{\texttt{#1}}
\newcommand{\datei}[1]{\texttt{#1}}
\newcommand{\minutes}[1]{\hfill{\normalsize\mdseries\color{ohmgrey}\faClock[regular]~#1 min}}

% ---------------------------------------------------------------- boxes
\tcbset{base/.style={enhanced, breakable, boxrule=0pt, arc=1.5pt, left=8pt, right=8pt,
        top=5pt, bottom=5pt, before skip=8pt, after skip=8pt}}
% \begin{todo}{1.3}: the block STUDENT 1.3, with its hint command at the end of the last line
\newcommand{\hintcmd}[1]{\unskip\nobreak\hfil\penalty50\hskip1em\null\nobreak\hfil
        \mbox{\small\color{ohmgrey}\faLightbulb[regular]~\texttt{python3 hint.py #1}}{\parfillskip=0pt\par}}
\newtcolorbox{todo}[1]{base, colback=ohmred!5, borderline west={2.5pt}{0pt}{ohmred},
        title={\faPen~STUDENT #1}, coltitle=ohmred, colbacktitle=ohmred!5, fonttitle=\bfseries,
        attach title to upper={\quad}, after upper={\hintcmd{#1}}}
\newtcolorbox{script}[1]{base, colback=black!3, borderline west={2.5pt}{0pt}{ohmgrey},
        title={\faBook~Lecture script, #1}, coltitle=ohmgrey, colbacktitle=black!3,
        fonttitle=\small\bfseries, toptitle=2pt, bottomtitle=0pt}
\newtcolorbox{note}{base, colback=ohmamber!8, borderline west={2.5pt}{0pt}{ohmamber}}
\newtcolorbox{checkpoint}{base, colback=ohmgreen!7, borderline west={2.5pt}{0pt}{ohmgreen},
        title={\faCheckCircle[regular]~Checkpoint}, coltitle=ohmgreen, colbacktitle=ohmgreen!7,
        fonttitle=\bfseries, attach title to upper={\quad}}

% \running{<launch arguments>}{<node arguments>}: what terminals 1 and 2 run during a task
\newcommand{\running}[2]{\begin{tcolorbox}[base, colback=black!3, borderline west={2.5pt}{0pt}{ohmgrey},
        title={\faTerminal~Running}, coltitle=ohmgrey, colbacktitle=black!3, fonttitle=\small\bfseries,
        toptitle=2pt, bottomtitle=0pt, fontupper=\small]
  \begin{tabularx}{\linewidth}{@{}>{\color{ohmgrey}}l >{\raggedright\arraybackslash}X@{}}
    terminal 1 & \texttt{ros2 launch ohm\_polygon\_map box.launch.py#1} \\
    terminal 2 & \texttt{ros2 run ohm\_polygon\_map polygon\_map\_node.py -\kern0pt-ros-args #2}
  \end{tabularx}
\end{tcolorbox}}
% \expected[<note>]{<name>}: RViz (figures/rviz_<name>.tikz) and terminal 2 (figures/term_<name>.txt)
\newtcolorbox{seebox}{base, colback=ohmgreen!3, borderline west={2.5pt}{0pt}{ohmgreen!60},
        title={\faEye[regular]~What you should see}, coltitle=ohmgreen!80!black, colbacktitle=ohmgreen!3,
        fonttitle=\small\bfseries, toptitle=2pt, bottomtitle=0pt, breakable=false}
\newcommand{\expected}[2][]{\begin{seebox}
  \begin{minipage}[c]{0.36\linewidth}\resizebox{\linewidth}{!}{\input{figures/rviz_#2.tikz}}\end{minipage}%
  \hfill
  \begin{minipage}[c]{0.615\linewidth}
    {\footnotesize\color{ohmgrey}#1terminal 2:}
    \lstinputlisting[basicstyle=\ttfamily\scriptsize, aboveskip=2pt, belowskip=0pt, xleftmargin=4pt,
                     framexleftmargin=3pt, breakindent=2em]{figures/term_#2.txt}
  \end{minipage}
\end{seebox}}

\setlist[enumerate,1]{leftmargin=*,itemsep=3pt,topsep=2pt}
\setlist[itemize,1]{leftmargin=*,itemsep=2pt,topsep=2pt}

\begin{document}
\sloppy

% ---------------------------------------------------------------- title
\begin{tcolorbox}[enhanced, colback=ohmred, colframe=ohmred, arc=0pt, left=14pt, right=14pt,
                  top=12pt, bottom=12pt, fontupper=\color{white}]
  {\small Intelligent Robotics (ASY2) · Lab 2}\\[6pt]
  {\LARGE\bfseries Polygon map}\\[4pt]
  {\large Least squares and RANSAC on a laser scan}
\end{tcolorbox}

\noindent
\begin{tabular}{@{}>{\color{ohmgrey}}l l@{\hspace{2.4em}}>{\color{ohmgrey}}l l@{}}
  \faClock[regular]~Duration & 180 minutes & \faUsers~Work in & pairs \\
  \faBook~Script & Cycle 2, section 2.2.3 & \faFolderOpen[regular]~Files & \datei{ohm-polygon-map/student/}
\end{tabular}

\section{Overview}
A robot stands in a room. We know two things about the room: it is a rectangle, and there is one round
object in it. From a single laser scan you build a map of the room:
\begin{itemize}
  \item the four walls as straight lines,
  \item the four corners as the intersections of these lines,
  \item the round object as a circle.
\end{itemize}
One corner is hidden behind the object (\cref{fig:hall}). No beam reaches it, but your map will contain
it anyway.

\begin{figure}[ht]
  \centering
  % TikZ figure: the hall `box` from above, in hall metres from the inner lower left corner.
% Robot at (2.0, 1.5), heading 20 deg; circle at (6.0, 3.8), r = 0.8 (config/box.json minus the 0.5 m wall).
% The shadow is bounded by the two tangents from the laser to the circle (19.9 and 39.9 deg).
\begin{tikzpicture}[font=\small, scale=0.95]
  \fill[black!55] (-0.5,-0.5) rectangle (9.5,6.5);
  \fill[white] (0,0) rectangle (9,6);
  \fill[black!7] (6.273,3.047) -- (9,4.033) -- (9,6) -- (7.380,6) -- (5.486,4.414) -- cycle;
  \draw[black!40, dashed] (2,1.5) -- (9,4.033);
  \draw[black!40, dashed] (2,1.5) -- (7.380,6);
  \fill[black!55] (6,3.8) circle (0.8);
  \node[white, align=center, font=\footnotesize] at (6,3.8) {round\\object};
  \node[black!55, font=\footnotesize] at (8.1,5.2) {shadow};
  % beam 250: straight down onto the lower wall
  \draw[ohmamber, line width=0.9pt] (2,1.5) -- (2,0);
  \node[ohmamber, font=\footnotesize, anchor=west] at (2.05,0.45) {beam 250};
  % robot and laser frame
  \draw[thick, fill=ohmred!15] (2,1.5) circle (0.23);
  \draw[black!40] (2,1.5) -- (3.7,1.5);
  \draw[black!60] (3.45,1.5) arc (0:20:1.45) node[pos=0.4, right, font=\footnotesize] {\(20^\circ\)};
  \begin{scope}[shift={(2,1.5)}, rotate=20]
    \draw[-{Stealth[length=2.2mm]}, thick, ohmred] (0,0) -- (1.0,0) node[above] {\(x\)};
    \draw[-{Stealth[length=2.2mm]}, thick, ohmgreen] (0,0) -- (0,1.0) node[above] {\(y\)};
  \end{scope}
  \node[font=\footnotesize, anchor=east] at (1.72,1.35) {laser};
  % corners
  \foreach \c in {(0,0), (9,0), (0,6)} \fill[ohmred] \c circle (0.11);
  \draw[ohmred, thick] (9,6) circle (0.17);
  \node[ohmred, font=\footnotesize, anchor=north east] at (8.85,5.95) {hidden corner};
  % dimensions
  \draw[{Stealth[length=2mm]}-{Stealth[length=2mm]}] (0,-0.95) -- (9,-0.95) node[midway, below] {9\,m};
  \draw[{Stealth[length=2mm]}-{Stealth[length=2mm]}] (9.95,0) -- (9.95,6) node[midway, right] {6\,m};
\end{tikzpicture}

  \caption{The hall \thema{box} from above. The laser looks along \(x\); all points, lines and circles
    in this lab are in its frame.}
  \label{fig:hall}
\end{figure}

\begin{table}[ht]\centering
\caption{Schedule.}
\label{tab:schedule}
\begin{tabular}{@{}rll@{}}
\toprule
min & task & done when \\
\midrule
0--20 & \cref{sec:task0} The scan and its noise & \cref{tab:noise} is filled in \\
20--60 & \cref{sec:task1} RANSAC: one line & a red line lies on a wall \\
60--85 & \cref{sec:task2} Least squares & the refined line fits better \\
85--95 & break & \\
95--120 & \cref{sec:task3} Four lines & four walls in four colours \\
120--140 & \cref{sec:task4} Corners & a red rectangle, \(9 \times 6\)\,m \\
140--165 & \cref{sec:task5} The circle & a red circle on the object \\
165--180 & \cref{sec:extra} Extra tasks & \\
\bottomrule
\end{tabular}
\end{table}

\paragraph{Notation} as in the script: a point is \((\hat x_i, \hat y_i)\), a line is
\(y = \tilde m\,x + \tilde t\), a circle has the centre \((\tilde x_0, \tilde y_0)\) and the radius
\(\tilde r\). In the code: points \thema{(x, y)}, lines \thema{(m, t)}, circles \thema{(x0, y0, r)}, all
in metres in the frame of the laser. The equations you need are repeated in the tasks, with the numbers
they have in the script.

\newpage
\section{Getting started}
\paragraph{Your own copy.} Fork \datei{ohm-polygon-map} on GitHub (button \emph{Fork} at the top right),
then clone your fork into the \datei{src/} folder of your workspace. The simulator is still there from
lab 1; update it:
\begin{lstlisting}
git clone git@github.com:<your-name>/ohm-polygon-map.git
\end{lstlisting}
\begin{lstlisting}
git -C ir-mecanum-lab pull
\end{lstlisting}
Commit after every task and push, so nothing is lost at the end of the session.

\paragraph{Build.} Build once, then source the workspace in every new terminal:
\begin{lstlisting}
colcon build --symlink-install --packages-select mecanum_lab_interfaces mecanum_lab ohm_polygon_map
\end{lstlisting}
\begin{lstlisting}
source install/setup.bash
\end{lstlisting}
You do not need to build again after a change to your code. Just start the node again.

\paragraph{Simulator.} One launch file starts the simulator in the hall \thema{box} and RViz. The robot
stands still; the keys do not move it.
\begin{lstlisting}
ros2 launch ohm_polygon_map box.launch.py
\end{lstlisting}

\paragraph{Your files.} \datei{student/line\_fitting.py} (tasks 0 to 4) and
\datei{student/circle\_fitting.py} (task 5), plain Python without ROS. In every function the checks
(too few points, \thema{None}, parallel lines, \dots) are already there. You write only the few lines
in the middle, marked
\begin{lstlisting}
# the slope from the two points, then the offset from one point that lies on the line
# TODO STUDENT 1.1
raise NotImplementedError("STUDENT 1.1")
\end{lstlisting}
The comment above it says in words what goes there; the equation is in the task of this handout, and
the \thema{return} below shows the names to use. Replace the \thema{raise} line by your code. Use only Python
and \thema{math}, no \thema{numpy}. The helpers you call are given in
\datei{ohm\_polygon\_map/geometry.py}, the drawing in RViz in \datei{ohm\_polygon\_map/rviz\_draw.py}.
You do not change these files.

\paragraph{The node.} \datei{student/polygon\_map\_node.py} is given. It subscribes to the scan, calls
your functions and draws the result in RViz. The parameter \thema{stage} says how far it goes, from 0 (the
noise) to 5 (the circle). It runs in a second terminal:
\begin{lstlisting}
ros2 run ohm_polygon_map polygon_map_node.py --ros-args -p stage:=1
\end{lstlisting}
Every 2\,s it takes the newest scan and prints what it found. If a function is not written yet, the node
stops with \thema{NotImplementedError: STUDENT 1.1}: that is the next one to write.

\paragraph{Self check.} The tests try your functions on made-up points, one stage at a time:
\begin{lstlisting}
python3 -m pytest test/test_fitting.py -q -k stage1
\end{lstlisting}

\paragraph{Three terminals.} All of them in the workspace, sourced.
\begin{center}\small
\begin{tabular}{@{}>{\color{ohmgrey}}l l l@{}}
  terminal 1 & simulator and RViz & stays open the whole session \\
  terminal 2 & your node           & Ctrl-C and start again after every change \\
  terminal 3 & tests, hints, \thema{ros2 topic} & \\
\end{tabular}
\end{center}
Every task starts with a box \emph{Running}: what terminals 1 and 2 run. It ends with \emph{What you
should see}: RViz and terminal 2 when the task is done. RViz shows only the scan and what your node
draws, not the walls of the hall. The numbers change a little from scan to scan.

\paragraph{Hints.} Stuck? Every block has three hints: the idea, how to write it in Python, and at last
the code. Each call opens the next one, in the folder \datei{ohm-polygon-map}:
\begin{lstlisting}
python3 hint.py 1.3
\end{lstlisting}
\thema{python3 hint.py} alone shows the hints you have opened. Try on your own first: the hints notice
when you ask too early.

\begin{figure}[!hb]
  \centering
  \resizebox{0.8\linewidth}{!}{% TikZ figure: from the scan to the map. Red boxes are the functions you write (task number in the
% corner), grey boxes are given. One evaluation every 2 s; tau comes once from task 0.
\begin{tikzpicture}[
  font=\small,
  data/.style={rectangle, rounded corners=2pt, draw=black!45, fill=black!3, minimum height=0.8cm,
               minimum width=2.6cm, align=center, inner sep=3pt},
  mine/.style={rectangle, draw=ohmred, line width=0.8pt, fill=ohmred!7, minimum height=0.8cm,
               minimum width=3.0cm, align=center, inner sep=3pt},
  given/.style={rectangle, draw=black!45, fill=white, minimum height=0.8cm, minimum width=3.0cm,
                align=center, inner sep=3pt},
  flow/.style={-{Stealth[length=2.2mm]}, line width=0.7pt, draw=black!70},
  tag/.style={font=\scriptsize\bfseries, text=ohmred, anchor=north east, inner sep=1.5pt},
]
\node[data]  (scan)   at (0, 0)     {\texttt{/muster/scan}\\[-1pt]{\scriptsize LaserScan, 20\,Hz}};
\node[given] (pts)    at (4.0, 0)   {\texttt{scan\_to\_points}};
\node[data]  (points) at (8.0, 0)   {points \((x_i, y_i)\)};
\node[mine]  (noise)  at (4.0, 1.6) {\texttt{mean\_and\_std}\\[-1pt]{\scriptsize one beam, 100 scans}};
\node[data]  (tau)    at (8.0, 1.6) {\(\tau = 3\sigma\)};

\node[mine, minimum width=7.6cm, minimum height=2.2cm] (find) at (8.0, -2.4) {};
\node[font=\scriptsize\ttfamily, anchor=north west] at (find.north west) {find\_lines: 4 times};
\node[mine, fill=white] (ransac) at (6.1, -2.55) {\texttt{ransac\_line}\\[-1pt]{\scriptsize\texttt{line\_through}, \texttt{num\_iterations}}};
\node[mine, fill=white, minimum width=2.4cm] (ols) at (10.3, -2.55) {\texttt{fit\_line}\\[-1pt]{\scriptsize on the inliers}};

\node[mine]  (cut)    at (6.0, -5.0) {\texttt{intersect}\\[-1pt]{\scriptsize in \texttt{room\_corners}}};
\node[data]  (room)   at (6.0, -6.6) {4 corners, the room polygon};
\node[mine]  (circ)   at (11.0, -5.0) {\texttt{ransac\_circle}\\[-1pt]{\scriptsize \texttt{circle\_distance}}};
\node[data]  (circle) at (11.0, -6.6) {circle \((x_0, y_0, r)\)};
\node[data, minimum width=7.0cm] (rviz) at (8.5, -8.0) {RViz: \texttt{/muster/polygon\_map/...}};

\node[tag] at (noise.north east) {0};
\node[tag] at (ransac.north east) {1};
\node[tag] at (ols.north east) {2};
\node[tag] at (find.north east) {3};
\node[tag] at (cut.north east) {4};
\node[tag] at (circ.north east) {5};

\draw[flow] (scan) -- (pts);
\draw[flow] (pts) -- (points);
\draw[flow] (scan.north) |- (noise.west);
\draw[flow] (noise) -- (tau);
\draw[flow] (points) -- (find.north -| points);
\draw[flow] (tau.east) -- ++(4.4,0) |- (find.east);
\draw[flow] (ransac) -- (ols);
\draw[flow] (find.south -| cut) -- node[left, font=\scriptsize] {4 lines} (cut);
\draw[flow] (find.south -| circ) -- node[right, font=\scriptsize] {points no line explains} (circ);
\draw[flow] (cut) -- (room);
\draw[flow] (circ) -- (circle);
\draw[flow] (room.south) -- (room.south |- rviz.north);
\draw[flow] (circle.south) -- (circle.south |- rviz.north);
\end{tikzpicture}
}
  \caption{From the scan to the map. Red: you write it (the number is the task). Grey: given.}
  \label{fig:pipeline}
\end{figure}

\clearpage
\section{Task 0: the scan and its noise\minutes{20}}
\label{sec:task0}
\running{}{-p stage:=0}

\begin{enumerate}
  \item In terminal 3, look at one scan:
\begin{lstlisting}
ros2 topic echo /muster/scan --once
\end{lstlisting}
        Fill in the first rows of \cref{tab:noise}. How many beams are there? In which direction does
        beam 0 point, in which beam 90?
  \item Read \thema{scan\_to\_points} in \datei{geometry.py}. Which beams does it drop, and why?
\end{enumerate}

A wall that does not move should give the same range in every scan. It does not: the laser is noisy. You
measure how noisy, because RANSAC needs a threshold \(\tau\) for ``close to the line''.

\begin{script}{section 2.2.3, the one-dimensional case}
The least-squares estimate of a constant from the readings \(\hat x_1 \dots \hat x_n\) is their mean:
\begin{equation}
  \tilde c = \frac 1n \sum_{i=1}^n \hat x_i .
  \tag{2.18}
\end{equation}
\end{script}

\begin{todo}{0.1}
\datei{line\_fitting.py}, \thema{mean\_and\_std(values)}: the mean (2.18) and the sample standard
deviation \(\sigma = \sqrt{\frac{1}{n-1}\sum_i (\hat x_i - \tilde c)^2}\).
\end{todo}

Run the node with \thema{stage:=0}. It collects beam 250 over 100 scans (5\,s), which hits the lower wall
at a right angle (\cref{fig:hall}). It prints the mean, \(\sigma\) and \(\tau = 3\sigma\). Repeat with
another beam, e.g. \thema{-p noise\_beam:=90}.

\begin{table}[ht]\centering
\renewcommand{\arraystretch}{1.8}
\caption{The scan and its noise. Fill in.}
\label{tab:noise}
\begin{tabular}{@{}>{\raggedright\arraybackslash}p{6.2cm}p{3.2cm}p{3.2cm}@{}}
\toprule
& beam 250 & beam 90 \\
\midrule
\thema{angle\_increment}, number of beams & \multicolumn{2}{l}{} \\
\thema{range\_max}, a beam without echo & \multicolumn{2}{l}{} \\
mean range \(\tilde c\) & & \\
\(\sigma\) & & \\
\(\tau = 3\sigma\) & & \\
\bottomrule
\end{tabular}
\end{table}

\expected[tick \thema{raw scan} in the list of displays of RViz; ]{stage0}

\begin{checkpoint}
\(\sigma\) is about the same for both beams. Why \(3\sigma\)? What share of the points on a wall lies
closer than \(\tau\) to it?
\end{checkpoint}

\newpage
\section{Task 1: RANSAC, one line\minutes{40}}
\label{sec:task1}
\running{}{-p stage:=1}

A scan of the room has points on four walls and on the object. A least-squares line through all of them
fits nothing. RANSAC turns the question around: draw two points, put a line through them, and count how
many points agree with it.

\begin{script}{section 2.2.3, robust fitting: RANSAC}
Repeat \(k\) times: \emph{sample} the minimum number of points (two for a line), \emph{fit} the model,
\emph{score} it by its inliers, the points closer than \(\tau\). \emph{Keep} the model with the most
inliers. To draw at least one all-inlier sample with probability \(p\), with inlier ratio \(w\) and sample
size \(s\):
\begin{equation}
  k = \frac{\log(1 - p)}{\log\left(1 - w^s\right)} .
  \tag{2.30}
\end{equation}
\end{script}

\begin{todo}{1.1}
\thema{line\_through(p, q)}: the line \(y = m x + t\) through two points,
\(m = (y_2 - y_1)/(x_2 - x_1)\), \(t = y_1 - m\,x_1\). Two points above each other are already caught.
\end{todo}

\begin{todo}{1.2}
\thema{num\_iterations(w, s, p)} with (2.30), rounded up (\thema{math.ceil}). Check with the exercise of
the script: \(w = 0.8\) gives 5, \(w = 0.5\) gives 17.
\end{todo}

\begin{note}
Given in \datei{geometry.py}: \thema{line\_distance(point, line)} is the distance perpendicular to the line,
\(d = |m x - y + t| / \sqrt{1 + m^2}\), and \thema{split\_line(points, line, tau)} returns
\thema{(inliers, outliers)}. \thema{rng.sample(points, 2)} draws two different points.
\end{note}

\begin{todo}{1.3}
\thema{ransac\_line(points, tau, iterations, rng)}: the loop already draws two points and puts a line
through them (\emph{sample}, \emph{fit}). You write \emph{score} and \emph{keep}: the inliers of the line,
and if there are more than for the best line so far, this line becomes the best line.
\end{todo}

Run the node with \thema{stage:=1}. RViz shows the scan in grey, the inliers and the line in red. The node
uses \(w = 0.3\) and \(p = 0.99\).
\begin{itemize}
  \item Which wall does RANSAC find? Why that one?
  \item Run it with \thema{-p seed:=2} up to \thema{seed:=5}. Does the result change?
  \item Set \thema{-p w\_line:=0.9}. How many iterations is that? What happens?
\end{itemize}

\expected[]{stage1}

\begin{checkpoint}
\thema{pytest -k stage1} passes. In RViz the red line lies on a wall.
\end{checkpoint}

\newpage
\section{Task 2: least squares\minutes{25}}
\label{sec:task2}
\running{}{-p stage:=2}

The RANSAC line goes through two noisy points. All its inliers together say more: fit a line to them by
least squares.

\begin{script}{section 2.2.3, the two-dimensional case}
The line \(y = \tilde m x + \tilde t\) that minimises the squared error in \(y\),
\(\epsilon = \sum_i \big(\hat y_i - (\tilde m\,\hat x_i + \tilde t)\big)^2\), has
\begin{equation}
  \tilde m = \frac{\sum\limits_{i=1}^n (\hat x_i - \overline{\hat x})(\hat y_i - \overline{\hat y})}
                  {\sum\limits_{i=1}^n (\hat x_i - \overline{\hat x})^2},
  \tag{2.23}
\end{equation}
\begin{equation}
  \tilde t = \overline{\hat y} - \tilde m\,\overline{\hat x} .
  \tag{2.24}
\end{equation}
\end{script}

\begin{todo}{2.1}
\thema{fit\_line(points)}: \(\tilde m\) with (2.23) and \(\tilde t\) with (2.24). The means
\(\overline{\hat x}\), \(\overline{\hat y}\) are already computed.
\end{todo}

\thema{pytest -k stage2} also tries the exercise of the script: the points \((0, 1.1)\), \((1, 2.9)\),
\((2, 5.2)\), \((3, 6.8)\) give \(\tilde m = 1.94\) and \(\tilde t = 1.09\).

The node fits your line to the inliers of the RANSAC line and then takes the inliers of the new line.
Run it with \thema{stage:=2}: grey is the RANSAC line, red the least-squares line.
\begin{itemize}
  \item Compare \(m\) and \(t\) of both lines over a few evaluations. Which one changes less?
  \item Has the number of inliers changed?
\end{itemize}

\expected[]{stage2}

\begin{checkpoint}
\thema{pytest -k stage2} passes. The least-squares line is steadier from scan to scan than the RANSAC line.
\end{checkpoint}

\newpage
\section{Task 3: four lines\minutes{25}}
\label{sec:task3}
\running{}{-p stage:=3}

Each wall is an outlier for the other three walls. So after one wall is found, its inliers are removed,
and RANSAC runs again on the rest: \emph{sequential RANSAC}.

\begin{todo}{3.1}
\thema{find\_lines(points, n, tau, iterations, min\_inliers, rng)} repeats \(n\) times:
\begin{enumerate}
  \item \thema{ransac\_line} on the remaining points (given),
  \item stop if it has fewer than \thema{min\_inliers} (given),
  \item \thema{fit\_line} on its inliers,
  \item \thema{split\_line} of the remaining points with this line: the inliers belong to the wall,
        the outliers are the new remaining points,
  \item append \thema{(line, inliers)} to \thema{lines}.
\end{enumerate}
You write steps 3 to 5, three lines.
\end{todo}

Run the node with \thema{stage:=3}. Each wall gets its own colour. Fill in \cref{tab:rounds}: \(N\) is the
number of points at the start of the round (round 1: all points, the node prints them; then subtract the
inliers of the round before), \(n\) the inliers of the line found, \(w = n/N\), and \(k\) the iterations
(2.30) would ask for with this \(w\).

\begin{table}[ht]\centering
\renewcommand{\arraystretch}{1.8}
\caption{Sequential RANSAC, round by round. Fill in.}
\label{tab:rounds}
\begin{tabular}{@{}cp{2.6cm}p{2.0cm}p{2.0cm}p{2.0cm}p{2.0cm}@{}}
\toprule
round & wall & \(N\) & \(n\) & \(w = n/N\) & \(k\) \\
\midrule
1 & & & & & \\
2 & & & & & \\
3 & & & & & \\
4 & & & & & \\
\bottomrule
\end{tabular}
\end{table}

\begin{itemize}
  \item The node uses \(k = 49\) (\(w = 0.3\)) in every round. Too many, too few?
  \item Why does one wall have far fewer points than the others? (\cref{fig:hall})
  \item What would happen without \thema{min\_inliers}, if you asked for five lines?
\end{itemize}

\expected[]{stage3}

\begin{checkpoint}
\thema{pytest -k stage3} passes. Four walls in four colours, about 20 points left over.
\end{checkpoint}

\newpage
\section{Task 4: corners\minutes{20}}
\label{sec:task4}
\running{}{-p stage:=4}

Two walls meet in a corner: where both line equations hold,
\begin{equation}
  m_1 x + t_1 = m_2 x + t_2
  \quad\Rightarrow\quad
  x = \frac{t_2 - t_1}{m_1 - m_2}, \qquad y = m_1 x + t_1 .
  \label{eq:intersect}
\end{equation}
Four lines make six pairs. Two of them are opposite walls: parallel, no corner.

\begin{todo}{4.1}
\thema{intersect(l1, l2, min\_angle)}: \(x\) and \(y\) with \cref{eq:intersect}. Parallel lines (an
angle below \thema{min\_angle}) are already caught.
\end{todo}

The given \thema{room\_corners(lines, intersect)} calls your \thema{intersect} on all pairs, drops the \thema{None}s and orders the corners
around their centre, so that they form a polygon. Run the node with \thema{stage:=4}. RViz draws the room
polygon in red with the length of each side; the node also prints the angles at the corners.

\begin{table}[ht]\centering
\renewcommand{\arraystretch}{1.8}
\caption{The room. Fill in.}
\label{tab:room}
\begin{tabular}{@{}p{3.6cm}p{2.2cm}p{2.2cm}p{2.2cm}p{2.2cm}@{}}
\toprule
& side 1 & side 2 & side 3 & side 4 \\
\midrule
length in m & & & & \\
angle at its first corner & & & & \\
\bottomrule
\end{tabular}
\end{table}

\begin{itemize}
  \item Which corner was never measured? How far is it from the nearest scan point?
  \item Why is \thema{min\_angle} needed, and not \thema{m1 == m2}?
\end{itemize}

\expected[]{stage4}

\begin{checkpoint}
\thema{pytest -k stage4} passes. The sides are \(9\) and \(6\,\mathrm{m}\) within a few centimetres, the
angles \(90^\circ\) within a few tenths of a degree.
\end{checkpoint}

\newpage
\section{Task 5: the circle\minutes{25}}
\label{sec:task5}
\running{}{-p stage:=5}

After the four walls, about 20 points are left: the object. RANSAC works for a circle as it does for a
line. Only three things change.

\begin{script}{section 2.2.3, example: RANSAC for circles}
\emph{Sample}: \(s = 3\) points; three points not on a line define one circle.
\emph{Score}: point \(i\) is an inlier if
\(\big|\sqrt{(\hat x_i - \tilde x_0)^2 + (\hat y_i - \tilde y_0)^2} - \tilde r\big| < \tau\).
\emph{Refit}: the least-squares circle (2.37), (2.38) on all inliers. Applied to three points, the same
formulas return the circle through them.
\end{script}

\begin{note}
Given in \datei{geometry.py}: \thema{fit\_circle(points)} is (2.37) and (2.38); it returns \thema{None} for
points on a line. Your file this time: \datei{student/circle\_fitting.py}.
\end{note}

\begin{todo}{5.1}
\thema{circle\_distance(point, circle)}: the distance of the \emph{score} in the script box.
\end{todo}

\begin{todo}{5.2}
\thema{ransac\_circle(points, tau, iterations, rng, r\_max)}: as in \thema{ransac\_line}, the loop already
draws three points and fits the circle through them; circles larger than \thema{r\_max} are skipped
(three points almost on a line give a huge circle). You write \emph{score} and \emph{keep}, with
\thema{circle\_distance}. The refit on all inliers at the end is given.
\end{todo}

Run the node with \thema{stage:=5}. The node uses \(w = 0.5\), so \(k = 35\).
\begin{itemize}
  \item Centre and radius of the circle? Compare with \datei{config/box.json} (hall coordinates, so it
        does not match directly: why?).
  \item How many iterations would (2.30) need on the \emph{whole} scan, \(w = 20/360\), \(s = 3\)? What
        does that say about the order of tasks 3 and 5?
\end{itemize}

\expected[]{stage5}

\begin{checkpoint}
\thema{pytest -k stage5} passes. A red circle lies on the object, the radius is \(0.8\,\mathrm{m}\) within a
few centimetres.
\end{checkpoint}

\newpage
\section{Extra tasks}
\label{sec:extra}

\subsection*{Z.1: a vertical wall}
Stop the simulator and start it with the robot turned to \(0^\circ\):
\running{ start:=straight}{-p stage:=4 \mbox{\textcolor{ohmgrey}{[-p hesse:=true]}}}
Run the node with \thema{stage:=4}. What happens to the map? Look at \(m\) of the walls that run along
\(y\). The script names the problem: \(y = m x + t\) measures the error only in \(y\), and it cannot
describe a vertical line. Robots use the Hesse normal form \(x\cos\alpha + y\sin\alpha = r\) instead.

\begin{script}{section 2.2.3, orthogonal least squares and the Hesse normal form}
\begin{equation}
  q_x = \sum_{i=1}^n (\hat x_i - \overline{\hat x})^2, \qquad
  q_y = \sum_{i=1}^n (\hat y_i - \overline{\hat y})^2, \qquad
  q_{xy} = \sum_{i=1}^n (\hat x_i - \overline{\hat x})(\hat y_i - \overline{\hat y}),
  \tag{2.25}
\end{equation}
\begin{equation}
  \alpha = \tfrac12 \operatorname{atan2}\!\left(-2\,q_{xy},\ q_y - q_x\right), \qquad
  r = \overline{\hat x}\cos\alpha + \overline{\hat y}\sin\alpha .
  \tag{2.29}
\end{equation}
If \(r\) comes out negative, add \(\pi\) to \(\alpha\) and flip the sign of \(r\).
\end{script}

\begin{todo}{Z.1}
\thema{fit\_line\_hesse(points)} in \datei{line\_fitting.py}: \(q_x, q_y, q_{xy}\) with (2.25), then \(\alpha\)
and \(r\) with (2.29). The means and the sign of \(r\) are already handled. Test with \thema{pytest -k extra}, then run the node
with \thema{-p stage:=4 -p hesse:=true}: tasks 3 and 4 now use your function instead of \thema{fit\_line}.
\end{todo}

\begin{seebox}
  \centering
  \begin{minipage}[t]{0.4\linewidth}\centering
    \resizebox{\linewidth}{!}{% made by tools/make_rviz_figures.py -- do not edit
\begin{tikzpicture}[x=0.62cm, y=0.62cm, font=\scriptsize]
  \fill[color={rgb,255:red,33;green,36;blue,43}] (-0.4,-0.4) rectangle (10.4,7.4);
  \clip (-0.4,-0.4) rectangle (10.4,7.4);
  \draw[color={rgb,255:red,68;green,73;blue,84}, line width=0.53pt, shift={(0,0.5)}] (-1,-3) grid (11,9);
  \draw[color={rgb,255:red,230;green,40;blue,40}, line width=1.2pt] (0.000,0.000) -- (0.600,0.000);
  \draw[color={rgb,255:red,40;green,200;blue,60}, line width=1.2pt] (0.000,0.000) -- (0.000,0.600);
  \draw[color={rgb,255:red,230;green,40;blue,40}, line width=1.2pt] (2.500,2.000) -- (3.100,2.000);
  \draw[color={rgb,255:red,40;green,200;blue,60}, line width=1.2pt] (2.500,2.000) -- (2.500,2.600);
  \fill[color={rgb,255:red,153;green,153;blue,153}] (9.489,1.975) rectangle ++(0.050,0.050);
  \fill[color={rgb,255:red,153;green,153;blue,153}] (9.490,2.097) rectangle ++(0.050,0.050);
  \fill[color={rgb,255:red,153;green,153;blue,153}] (9.480,2.220) rectangle ++(0.050,0.050);
  \fill[color={rgb,255:red,153;green,153;blue,153}] (9.499,2.343) rectangle ++(0.050,0.050);
  \fill[color={rgb,255:red,153;green,153;blue,153}] (9.464,2.464) rectangle ++(0.050,0.050);
  \fill[color={rgb,255:red,153;green,153;blue,153}] (9.486,2.588) rectangle ++(0.050,0.050);
  \fill[color={rgb,255:red,153;green,153;blue,153}] (9.457,2.709) rectangle ++(0.050,0.050);
  \fill[color={rgb,255:red,153;green,153;blue,153}] (9.481,2.835) rectangle ++(0.050,0.050);
  \fill[color={rgb,255:red,153;green,153;blue,153}] (9.478,2.959) rectangle ++(0.050,0.050);
  \fill[color={rgb,255:red,153;green,153;blue,153}] (9.488,3.086) rectangle ++(0.050,0.050);
  \fill[color={rgb,255:red,153;green,153;blue,153}] (9.492,3.212) rectangle ++(0.050,0.050);
  \fill[color={rgb,255:red,153;green,153;blue,153}] (9.502,3.341) rectangle ++(0.050,0.050);
  \fill[color={rgb,255:red,153;green,153;blue,153}] (9.480,3.464) rectangle ++(0.050,0.050);
  \fill[color={rgb,255:red,153;green,153;blue,153}] (9.484,3.593) rectangle ++(0.050,0.050);
  \fill[color={rgb,255:red,153;green,153;blue,153}] (9.481,3.722) rectangle ++(0.050,0.050);
  \fill[color={rgb,255:red,153;green,153;blue,153}] (9.432,3.839) rectangle ++(0.050,0.050);
  \fill[color={rgb,255:red,153;green,153;blue,153}] (9.469,3.980) rectangle ++(0.050,0.050);
  \fill[color={rgb,255:red,153;green,153;blue,153}] (9.471,4.114) rectangle ++(0.050,0.050);
  \fill[color={rgb,255:red,153;green,153;blue,153}] (9.456,4.243) rectangle ++(0.050,0.050);
  \fill[color={rgb,255:red,153;green,153;blue,153}] (9.484,4.388) rectangle ++(0.050,0.050);
  \fill[color={rgb,255:red,153;green,153;blue,153}] (6.639,3.491) rectangle ++(0.050,0.050);
  \fill[color={rgb,255:red,153;green,153;blue,153}] (6.387,3.477) rectangle ++(0.050,0.050);
  \fill[color={rgb,255:red,153;green,153;blue,153}] (6.258,3.504) rectangle ++(0.050,0.050);
  \fill[color={rgb,255:red,153;green,153;blue,153}] (6.174,3.545) rectangle ++(0.050,0.050);
  \fill[color={rgb,255:red,153;green,153;blue,153}] (6.085,3.582) rectangle ++(0.050,0.050);
  \fill[color={rgb,255:red,153;green,153;blue,153}] (6.016,3.626) rectangle ++(0.050,0.050);
  \fill[color={rgb,255:red,153;green,153;blue,153}] (5.956,3.673) rectangle ++(0.050,0.050);
  \fill[color={rgb,255:red,153;green,153;blue,153}] (5.894,3.717) rectangle ++(0.050,0.050);
  \fill[color={rgb,255:red,153;green,153;blue,153}] (5.841,3.765) rectangle ++(0.050,0.050);
  \fill[color={rgb,255:red,153;green,153;blue,153}] (5.787,3.811) rectangle ++(0.050,0.050);
  \fill[color={rgb,255:red,153;green,153;blue,153}] (5.782,3.884) rectangle ++(0.050,0.050);
  \fill[color={rgb,255:red,153;green,153;blue,153}] (5.733,3.932) rectangle ++(0.050,0.050);
  \fill[color={rgb,255:red,153;green,153;blue,153}] (5.725,4.006) rectangle ++(0.050,0.050);
  \fill[color={rgb,255:red,153;green,153;blue,153}] (5.701,4.070) rectangle ++(0.050,0.050);
  \fill[color={rgb,255:red,153;green,153;blue,153}] (5.654,4.119) rectangle ++(0.050,0.050);
  \fill[color={rgb,255:red,153;green,153;blue,153}] (5.681,4.220) rectangle ++(0.050,0.050);
  \fill[color={rgb,255:red,153;green,153;blue,153}] (5.686,4.308) rectangle ++(0.050,0.050);
  \fill[color={rgb,255:red,153;green,153;blue,153}] (5.686,4.395) rectangle ++(0.050,0.050);
  \fill[color={rgb,255:red,153;green,153;blue,153}] (5.706,4.499) rectangle ++(0.050,0.050);
  \fill[color={rgb,255:red,153;green,153;blue,153}] (5.754,4.630) rectangle ++(0.050,0.050);
  \fill[color={rgb,255:red,153;green,153;blue,153}] (7.838,6.475) rectangle ++(0.050,0.050);
  \fill[color={rgb,255:red,153;green,153;blue,153}] (7.655,6.478) rectangle ++(0.050,0.050);
  \fill[color={rgb,255:red,153;green,153;blue,153}] (7.490,6.490) rectangle ++(0.050,0.050);
  \fill[color={rgb,255:red,153;green,153;blue,153}] (7.282,6.458) rectangle ++(0.050,0.050);
  \fill[color={rgb,255:red,153;green,153;blue,153}] (7.126,6.467) rectangle ++(0.050,0.050);
  \fill[color={rgb,255:red,153;green,153;blue,153}] (6.969,6.469) rectangle ++(0.050,0.050);
  \fill[color={rgb,255:red,153;green,153;blue,153}] (6.825,6.480) rectangle ++(0.050,0.050);
  \fill[color={rgb,255:red,153;green,153;blue,153}] (6.680,6.484) rectangle ++(0.050,0.050);
  \fill[color={rgb,255:red,153;green,153;blue,153}] (6.543,6.493) rectangle ++(0.050,0.050);
  \fill[color={rgb,255:red,153;green,153;blue,153}] (6.395,6.484) rectangle ++(0.050,0.050);
  \fill[color={rgb,255:red,153;green,153;blue,153}] (6.253,6.477) rectangle ++(0.050,0.050);
  \fill[color={rgb,255:red,153;green,153;blue,153}] (6.125,6.483) rectangle ++(0.050,0.050);
  \fill[color={rgb,255:red,153;green,153;blue,153}] (5.987,6.470) rectangle ++(0.050,0.050);
  \fill[color={rgb,255:red,153;green,153;blue,153}] (5.870,6.480) rectangle ++(0.050,0.050);
  \fill[color={rgb,255:red,153;green,153;blue,153}] (5.743,6.473) rectangle ++(0.050,0.050);
  \fill[color={rgb,255:red,153;green,153;blue,153}] (5.618,6.464) rectangle ++(0.050,0.050);
  \fill[color={rgb,255:red,153;green,153;blue,153}] (5.509,6.473) rectangle ++(0.050,0.050);
  \fill[color={rgb,255:red,153;green,153;blue,153}] (5.386,6.458) rectangle ++(0.050,0.050);
  \fill[color={rgb,255:red,153;green,153;blue,153}] (5.281,6.465) rectangle ++(0.050,0.050);
  \fill[color={rgb,255:red,153;green,153;blue,153}] (5.168,6.456) rectangle ++(0.050,0.050);
  \fill[color={rgb,255:red,153;green,153;blue,153}] (5.085,6.496) rectangle ++(0.050,0.050);
  \fill[color={rgb,255:red,153;green,153;blue,153}] (4.973,6.481) rectangle ++(0.050,0.050);
  \fill[color={rgb,255:red,153;green,153;blue,153}] (4.863,6.467) rectangle ++(0.050,0.050);
  \fill[color={rgb,255:red,153;green,153;blue,153}] (4.778,6.494) rectangle ++(0.050,0.050);
  \fill[color={rgb,255:red,153;green,153;blue,153}] (4.667,6.470) rectangle ++(0.050,0.050);
  \fill[color={rgb,255:red,153;green,153;blue,153}] (4.567,6.461) rectangle ++(0.050,0.050);
  \fill[color={rgb,255:red,153;green,153;blue,153}] (4.476,6.469) rectangle ++(0.050,0.050);
  \fill[color={rgb,255:red,153;green,153;blue,153}] (4.383,6.470) rectangle ++(0.050,0.050);
  \fill[color={rgb,255:red,153;green,153;blue,153}] (4.284,6.452) rectangle ++(0.050,0.050);
  \fill[color={rgb,255:red,153;green,153;blue,153}] (4.202,6.475) rectangle ++(0.050,0.050);
  \fill[color={rgb,255:red,153;green,153;blue,153}] (4.105,6.454) rectangle ++(0.050,0.050);
  \fill[color={rgb,255:red,153;green,153;blue,153}] (4.028,6.486) rectangle ++(0.050,0.050);
  \fill[color={rgb,255:red,153;green,153;blue,153}] (3.937,6.473) rectangle ++(0.050,0.050);
  \fill[color={rgb,255:red,153;green,153;blue,153}] (3.854,6.485) rectangle ++(0.050,0.050);
  \fill[color={rgb,255:red,153;green,153;blue,153}] (3.770,6.490) rectangle ++(0.050,0.050);
  \fill[color={rgb,255:red,153;green,153;blue,153}] (3.680,6.472) rectangle ++(0.050,0.050);
  \fill[color={rgb,255:red,153;green,153;blue,153}] (3.595,6.467) rectangle ++(0.050,0.050);
  \fill[color={rgb,255:red,153;green,153;blue,153}] (3.507,6.445) rectangle ++(0.050,0.050);
  \fill[color={rgb,255:red,153;green,153;blue,153}] (3.429,6.462) rectangle ++(0.050,0.050);
  \fill[color={rgb,255:red,153;green,153;blue,153}] (3.352,6.488) rectangle ++(0.050,0.050);
  \fill[color={rgb,255:red,153;green,153;blue,153}] (3.269,6.476) rectangle ++(0.050,0.050);
  \fill[color={rgb,255:red,153;green,153;blue,153}] (3.187,6.472) rectangle ++(0.050,0.050);
  \fill[color={rgb,255:red,153;green,153;blue,153}] (3.109,6.486) rectangle ++(0.050,0.050);
  \fill[color={rgb,255:red,153;green,153;blue,153}] (3.030,6.496) rectangle ++(0.050,0.050);
  \fill[color={rgb,255:red,153;green,153;blue,153}] (2.946,6.457) rectangle ++(0.050,0.050);
  \fill[color={rgb,255:red,153;green,153;blue,153}] (2.869,6.484) rectangle ++(0.050,0.050);
  \fill[color={rgb,255:red,153;green,153;blue,153}] (2.789,6.471) rectangle ++(0.050,0.050);
  \fill[color={rgb,255:red,153;green,153;blue,153}] (2.711,6.478) rectangle ++(0.050,0.050);
  \fill[color={rgb,255:red,153;green,153;blue,153}] (2.633,6.506) rectangle ++(0.050,0.050);
  \fill[color={rgb,255:red,153;green,153;blue,153}] (2.554,6.478) rectangle ++(0.050,0.050);
  \fill[color={rgb,255:red,153;green,153;blue,153}] (2.475,6.460) rectangle ++(0.050,0.050);
  \fill[color={rgb,255:red,153;green,153;blue,153}] (2.397,6.454) rectangle ++(0.050,0.050);
  \fill[color={rgb,255:red,153;green,153;blue,153}] (2.318,6.462) rectangle ++(0.050,0.050);
  \fill[color={rgb,255:red,153;green,153;blue,153}] (2.238,6.501) rectangle ++(0.050,0.050);
  \fill[color={rgb,255:red,153;green,153;blue,153}] (2.163,6.442) rectangle ++(0.050,0.050);
  \fill[color={rgb,255:red,153;green,153;blue,153}] (2.082,6.467) rectangle ++(0.050,0.050);
  \fill[color={rgb,255:red,153;green,153;blue,153}] (2.000,6.497) rectangle ++(0.050,0.050);
  \fill[color={rgb,255:red,153;green,153;blue,153}] (1.922,6.478) rectangle ++(0.050,0.050);
  \fill[color={rgb,255:red,153;green,153;blue,153}] (1.844,6.468) rectangle ++(0.050,0.050);
  \fill[color={rgb,255:red,153;green,153;blue,153}] (1.766,6.450) rectangle ++(0.050,0.050);
  \fill[color={rgb,255:red,153;green,153;blue,153}] (1.681,6.478) rectangle ++(0.050,0.050);
  \fill[color={rgb,255:red,153;green,153;blue,153}] (1.602,6.466) rectangle ++(0.050,0.050);
  \fill[color={rgb,255:red,153;green,153;blue,153}] (1.520,6.467) rectangle ++(0.050,0.050);
  \fill[color={rgb,255:red,153;green,153;blue,153}] (1.430,6.502) rectangle ++(0.050,0.050);
  \fill[color={rgb,255:red,153;green,153;blue,153}] (1.355,6.468) rectangle ++(0.050,0.050);
  \fill[color={rgb,255:red,153;green,153;blue,153}] (1.272,6.466) rectangle ++(0.050,0.050);
  \fill[color={rgb,255:red,153;green,153;blue,153}] (1.175,6.507) rectangle ++(0.050,0.050);
  \fill[color={rgb,255:red,153;green,153;blue,153}] (1.099,6.477) rectangle ++(0.050,0.050);
  \fill[color={rgb,255:red,153;green,153;blue,153}] (1.015,6.468) rectangle ++(0.050,0.050);
  \fill[color={rgb,255:red,153;green,153;blue,153}] (0.927,6.471) rectangle ++(0.050,0.050);
  \fill[color={rgb,255:red,153;green,153;blue,153}] (0.835,6.482) rectangle ++(0.050,0.050);
  \fill[color={rgb,255:red,153;green,153;blue,153}] (0.739,6.496) rectangle ++(0.050,0.050);
  \fill[color={rgb,255:red,153;green,153;blue,153}] (0.655,6.481) rectangle ++(0.050,0.050);
  \fill[color={rgb,255:red,153;green,153;blue,153}] (0.566,6.472) rectangle ++(0.050,0.050);
  \fill[color={rgb,255:red,153;green,153;blue,153}] (0.473,6.470) rectangle ++(0.050,0.050);
  \fill[color={rgb,255:red,153;green,153;blue,153}] (0.477,6.260) rectangle ++(0.050,0.050);
  \fill[color={rgb,255:red,153;green,153;blue,153}] (0.478,6.070) rectangle ++(0.050,0.050);
  \fill[color={rgb,255:red,153;green,153;blue,153}] (0.458,5.933) rectangle ++(0.050,0.050);
  \fill[color={rgb,255:red,153;green,153;blue,153}] (0.475,5.737) rectangle ++(0.050,0.050);
  \fill[color={rgb,255:red,153;green,153;blue,153}] (0.463,5.604) rectangle ++(0.050,0.050);
  \fill[color={rgb,255:red,153;green,153;blue,153}] (0.476,5.437) rectangle ++(0.050,0.050);
  \fill[color={rgb,255:red,153;green,153;blue,153}] (0.489,5.281) rectangle ++(0.050,0.050);
  \fill[color={rgb,255:red,153;green,153;blue,153}] (0.465,5.191) rectangle ++(0.050,0.050);
  \fill[color={rgb,255:red,153;green,153;blue,153}] (0.485,5.040) rectangle ++(0.050,0.050);
  \fill[color={rgb,255:red,153;green,153;blue,153}] (0.477,4.938) rectangle ++(0.050,0.050);
  \fill[color={rgb,255:red,153;green,153;blue,153}] (0.466,4.844) rectangle ++(0.050,0.050);
  \fill[color={rgb,255:red,153;green,153;blue,153}] (0.478,4.724) rectangle ++(0.050,0.050);
  \fill[color={rgb,255:red,153;green,153;blue,153}] (0.485,4.616) rectangle ++(0.050,0.050);
  \fill[color={rgb,255:red,153;green,153;blue,153}] (0.473,4.537) rectangle ++(0.050,0.050);
  \fill[color={rgb,255:red,153;green,153;blue,153}] (0.459,4.464) rectangle ++(0.050,0.050);
  \fill[color={rgb,255:red,153;green,153;blue,153}] (0.465,4.371) rectangle ++(0.050,0.050);
  \fill[color={rgb,255:red,153;green,153;blue,153}] (0.499,4.248) rectangle ++(0.050,0.050);
  \fill[color={rgb,255:red,153;green,153;blue,153}] (0.460,4.213) rectangle ++(0.050,0.050);
  \fill[color={rgb,255:red,153;green,153;blue,153}] (0.490,4.104) rectangle ++(0.050,0.050);
  \fill[color={rgb,255:red,153;green,153;blue,153}] (0.472,4.049) rectangle ++(0.050,0.050);
  \fill[color={rgb,255:red,153;green,153;blue,153}] (0.481,3.969) rectangle ++(0.050,0.050);
  \fill[color={rgb,255:red,153;green,153;blue,153}] (0.504,3.878) rectangle ++(0.050,0.050);
  \fill[color={rgb,255:red,153;green,153;blue,153}] (0.479,3.836) rectangle ++(0.050,0.050);
  \fill[color={rgb,255:red,153;green,153;blue,153}] (0.480,3.772) rectangle ++(0.050,0.050);
  \fill[color={rgb,255:red,153;green,153;blue,153}] (0.467,3.720) rectangle ++(0.050,0.050);
  \fill[color={rgb,255:red,153;green,153;blue,153}] (0.476,3.652) rectangle ++(0.050,0.050);
  \fill[color={rgb,255:red,153;green,153;blue,153}] (0.472,3.597) rectangle ++(0.050,0.050);
  \fill[color={rgb,255:red,153;green,153;blue,153}] (0.472,3.540) rectangle ++(0.050,0.050);
  \fill[color={rgb,255:red,153;green,153;blue,153}] (0.464,3.491) rectangle ++(0.050,0.050);
  \fill[color={rgb,255:red,153;green,153;blue,153}] (0.489,3.418) rectangle ++(0.050,0.050);
  \fill[color={rgb,255:red,153;green,153;blue,153}] (0.468,3.380) rectangle ++(0.050,0.050);
  \fill[color={rgb,255:red,153;green,153;blue,153}] (0.490,3.314) rectangle ++(0.050,0.050);
  \fill[color={rgb,255:red,153;green,153;blue,153}] (0.507,3.253) rectangle ++(0.050,0.050);
  \fill[color={rgb,255:red,153;green,153;blue,153}] (0.471,3.227) rectangle ++(0.050,0.050);
  \fill[color={rgb,255:red,153;green,153;blue,153}] (0.461,3.185) rectangle ++(0.050,0.050);
  \fill[color={rgb,255:red,153;green,153;blue,153}] (0.489,3.122) rectangle ++(0.050,0.050);
  \fill[color={rgb,255:red,153;green,153;blue,153}] (0.469,3.087) rectangle ++(0.050,0.050);
  \fill[color={rgb,255:red,153;green,153;blue,153}] (0.472,3.040) rectangle ++(0.050,0.050);
  \fill[color={rgb,255:red,153;green,153;blue,153}] (0.454,3.005) rectangle ++(0.050,0.050);
  \fill[color={rgb,255:red,153;green,153;blue,153}] (0.491,2.942) rectangle ++(0.050,0.050);
  \fill[color={rgb,255:red,153;green,153;blue,153}] (0.471,2.910) rectangle ++(0.050,0.050);
  \fill[color={rgb,255:red,153;green,153;blue,153}] (0.471,2.867) rectangle ++(0.050,0.050);
  \fill[color={rgb,255:red,153;green,153;blue,153}] (0.448,2.836) rectangle ++(0.050,0.050);
  \fill[color={rgb,255:red,153;green,153;blue,153}] (0.459,2.789) rectangle ++(0.050,0.050);
  \fill[color={rgb,255:red,153;green,153;blue,153}] (0.468,2.745) rectangle ++(0.050,0.050);
  \fill[color={rgb,255:red,153;green,153;blue,153}] (0.467,2.706) rectangle ++(0.050,0.050);
  \fill[color={rgb,255:red,153;green,153;blue,153}] (0.471,2.665) rectangle ++(0.050,0.050);
  \fill[color={rgb,255:red,153;green,153;blue,153}] (0.465,2.628) rectangle ++(0.050,0.050);
  \fill[color={rgb,255:red,153;green,153;blue,153}] (0.462,2.590) rectangle ++(0.050,0.050);
  \fill[color={rgb,255:red,153;green,153;blue,153}] (0.518,2.536) rectangle ++(0.050,0.050);
  \fill[color={rgb,255:red,153;green,153;blue,153}] (0.458,2.515) rectangle ++(0.050,0.050);
  \fill[color={rgb,255:red,153;green,153;blue,153}] (0.475,2.474) rectangle ++(0.050,0.050);
  \fill[color={rgb,255:red,153;green,153;blue,153}] (0.473,2.437) rectangle ++(0.050,0.050);
  \fill[color={rgb,255:red,153;green,153;blue,153}] (0.490,2.397) rectangle ++(0.050,0.050);
  \fill[color={rgb,255:red,153;green,153;blue,153}] (0.467,2.365) rectangle ++(0.050,0.050);
  \fill[color={rgb,255:red,153;green,153;blue,153}] (0.457,2.331) rectangle ++(0.050,0.050);
  \fill[color={rgb,255:red,153;green,153;blue,153}] (0.478,2.291) rectangle ++(0.050,0.050);
  \fill[color={rgb,255:red,153;green,153;blue,153}] (0.465,2.258) rectangle ++(0.050,0.050);
  \fill[color={rgb,255:red,153;green,153;blue,153}] (0.470,2.221) rectangle ++(0.050,0.050);
  \fill[color={rgb,255:red,153;green,153;blue,153}] (0.444,2.188) rectangle ++(0.050,0.050);
  \fill[color={rgb,255:red,153;green,153;blue,153}] (0.474,2.150) rectangle ++(0.050,0.050);
  \fill[color={rgb,255:red,153;green,153;blue,153}] (0.483,2.114) rectangle ++(0.050,0.050);
  \fill[color={rgb,255:red,153;green,153;blue,153}] (0.452,2.081) rectangle ++(0.050,0.050);
  \fill[color={rgb,255:red,153;green,153;blue,153}] (0.467,2.045) rectangle ++(0.050,0.050);
  \fill[color={rgb,255:red,153;green,153;blue,153}] (0.478,2.010) rectangle ++(0.050,0.050);
  \fill[color={rgb,255:red,153;green,153;blue,153}] (0.458,1.975) rectangle ++(0.050,0.050);
  \fill[color={rgb,255:red,153;green,153;blue,153}] (0.482,1.940) rectangle ++(0.050,0.050);
  \fill[color={rgb,255:red,153;green,153;blue,153}] (0.469,1.905) rectangle ++(0.050,0.050);
  \fill[color={rgb,255:red,153;green,153;blue,153}] (0.478,1.870) rectangle ++(0.050,0.050);
  \fill[color={rgb,255:red,153;green,153;blue,153}] (0.488,1.836) rectangle ++(0.050,0.050);
  \fill[color={rgb,255:red,153;green,153;blue,153}] (0.466,1.799) rectangle ++(0.050,0.050);
  \fill[color={rgb,255:red,153;green,153;blue,153}] (0.492,1.767) rectangle ++(0.050,0.050);
  \fill[color={rgb,255:red,153;green,153;blue,153}] (0.466,1.728) rectangle ++(0.050,0.050);
  \fill[color={rgb,255:red,153;green,153;blue,153}] (0.426,1.687) rectangle ++(0.050,0.050);
  \fill[color={rgb,255:red,153;green,153;blue,153}] (0.473,1.658) rectangle ++(0.050,0.050);
  \fill[color={rgb,255:red,153;green,153;blue,153}] (0.481,1.623) rectangle ++(0.050,0.050);
  \fill[color={rgb,255:red,153;green,153;blue,153}] (0.455,1.582) rectangle ++(0.050,0.050);
  \fill[color={rgb,255:red,153;green,153;blue,153}] (0.464,1.548) rectangle ++(0.050,0.050);
  \fill[color={rgb,255:red,153;green,153;blue,153}] (0.473,1.513) rectangle ++(0.050,0.050);
  \fill[color={rgb,255:red,153;green,153;blue,153}] (0.458,1.472) rectangle ++(0.050,0.050);
  \fill[color={rgb,255:red,153;green,153;blue,153}] (0.500,1.446) rectangle ++(0.050,0.050);
  \fill[color={rgb,255:red,153;green,153;blue,153}] (0.483,1.404) rectangle ++(0.050,0.050);
  \fill[color={rgb,255:red,153;green,153;blue,153}] (0.491,1.369) rectangle ++(0.050,0.050);
  \fill[color={rgb,255:red,153;green,153;blue,153}] (0.479,1.326) rectangle ++(0.050,0.050);
  \fill[color={rgb,255:red,153;green,153;blue,153}] (0.464,1.282) rectangle ++(0.050,0.050);
  \fill[color={rgb,255:red,153;green,153;blue,153}] (0.483,1.250) rectangle ++(0.050,0.050);
  \fill[color={rgb,255:red,153;green,153;blue,153}] (0.470,1.205) rectangle ++(0.050,0.050);
  \fill[color={rgb,255:red,153;green,153;blue,153}] (0.483,1.170) rectangle ++(0.050,0.050);
  \fill[color={rgb,255:red,153;green,153;blue,153}] (0.456,1.118) rectangle ++(0.050,0.050);
  \fill[color={rgb,255:red,153;green,153;blue,153}] (0.447,1.072) rectangle ++(0.050,0.050);
  \fill[color={rgb,255:red,153;green,153;blue,153}] (0.459,1.035) rectangle ++(0.050,0.050);
  \fill[color={rgb,255:red,153;green,153;blue,153}] (0.468,0.996) rectangle ++(0.050,0.050);
  \fill[color={rgb,255:red,153;green,153;blue,153}] (0.477,0.957) rectangle ++(0.050,0.050);
  \fill[color={rgb,255:red,153;green,153;blue,153}] (0.492,0.920) rectangle ++(0.050,0.050);
  \fill[color={rgb,255:red,153;green,153;blue,153}] (0.472,0.864) rectangle ++(0.050,0.050);
  \fill[color={rgb,255:red,153;green,153;blue,153}] (0.470,0.817) rectangle ++(0.050,0.050);
  \fill[color={rgb,255:red,153;green,153;blue,153}] (0.488,0.781) rectangle ++(0.050,0.050);
  \fill[color={rgb,255:red,153;green,153;blue,153}] (0.482,0.730) rectangle ++(0.050,0.050);
  \fill[color={rgb,255:red,153;green,153;blue,153}] (0.460,0.667) rectangle ++(0.050,0.050);
  \fill[color={rgb,255:red,153;green,153;blue,153}] (0.462,0.617) rectangle ++(0.050,0.050);
  \fill[color={rgb,255:red,153;green,153;blue,153}] (0.485,0.581) rectangle ++(0.050,0.050);
  \fill[color={rgb,255:red,153;green,153;blue,153}] (0.469,0.517) rectangle ++(0.050,0.050);
  \fill[color={rgb,255:red,153;green,153;blue,153}] (0.470,0.464) rectangle ++(0.050,0.050);
  \fill[color={rgb,255:red,153;green,153;blue,153}] (0.537,0.461) rectangle ++(0.050,0.050);
  \fill[color={rgb,255:red,153;green,153;blue,153}] (0.602,0.458) rectangle ++(0.050,0.050);
  \fill[color={rgb,255:red,153;green,153;blue,153}] (0.681,0.470) rectangle ++(0.050,0.050);
  \fill[color={rgb,255:red,153;green,153;blue,153}] (0.768,0.491) rectangle ++(0.050,0.050);
  \fill[color={rgb,255:red,153;green,153;blue,153}] (0.799,0.466) rectangle ++(0.050,0.050);
  \fill[color={rgb,255:red,153;green,153;blue,153}] (0.891,0.498) rectangle ++(0.050,0.050);
  \fill[color={rgb,255:red,153;green,153;blue,153}] (0.938,0.491) rectangle ++(0.050,0.050);
  \fill[color={rgb,255:red,153;green,153;blue,153}] (0.978,0.478) rectangle ++(0.050,0.050);
  \fill[color={rgb,255:red,153;green,153;blue,153}] (1.038,0.487) rectangle ++(0.050,0.050);
  \fill[color={rgb,255:red,153;green,153;blue,153}] (1.075,0.473) rectangle ++(0.050,0.050);
  \fill[color={rgb,255:red,153;green,153;blue,153}] (1.114,0.464) rectangle ++(0.050,0.050);
  \fill[color={rgb,255:red,153;green,153;blue,153}] (1.175,0.480) rectangle ++(0.050,0.050);
  \fill[color={rgb,255:red,153;green,153;blue,153}] (1.228,0.489) rectangle ++(0.050,0.050);
  \fill[color={rgb,255:red,153;green,153;blue,153}] (1.260,0.474) rectangle ++(0.050,0.050);
  \fill[color={rgb,255:red,153;green,153;blue,153}] (1.294,0.464) rectangle ++(0.050,0.050);
  \fill[color={rgb,255:red,153;green,153;blue,153}] (1.361,0.496) rectangle ++(0.050,0.050);
  \fill[color={rgb,255:red,153;green,153;blue,153}] (1.391,0.483) rectangle ++(0.050,0.050);
  \fill[color={rgb,255:red,153;green,153;blue,153}] (1.428,0.480) rectangle ++(0.050,0.050);
  \fill[color={rgb,255:red,153;green,153;blue,153}] (1.483,0.505) rectangle ++(0.050,0.050);
  \fill[color={rgb,255:red,153;green,153;blue,153}] (1.500,0.474) rectangle ++(0.050,0.050);
  \fill[color={rgb,255:red,153;green,153;blue,153}] (1.553,0.499) rectangle ++(0.050,0.050);
  \fill[color={rgb,255:red,153;green,153;blue,153}] (1.574,0.476) rectangle ++(0.050,0.050);
  \fill[color={rgb,255:red,153;green,153;blue,153}] (1.625,0.503) rectangle ++(0.050,0.050);
  \fill[color={rgb,255:red,153;green,153;blue,153}] (1.646,0.479) rectangle ++(0.050,0.050);
  \fill[color={rgb,255:red,153;green,153;blue,153}] (1.664,0.449) rectangle ++(0.050,0.050);
  \fill[color={rgb,255:red,153;green,153;blue,153}] (1.717,0.488) rectangle ++(0.050,0.050);
  \fill[color={rgb,255:red,153;green,153;blue,153}] (1.735,0.457) rectangle ++(0.050,0.050);
  \fill[color={rgb,255:red,153;green,153;blue,153}] (1.780,0.485) rectangle ++(0.050,0.050);
  \fill[color={rgb,255:red,153;green,153;blue,153}] (1.808,0.477) rectangle ++(0.050,0.050);
  \fill[color={rgb,255:red,153;green,153;blue,153}] (1.847,0.496) rectangle ++(0.050,0.050);
  \fill[color={rgb,255:red,153;green,153;blue,153}] (1.885,0.514) rectangle ++(0.050,0.050);
  \fill[color={rgb,255:red,153;green,153;blue,153}] (1.897,0.469) rectangle ++(0.050,0.050);
  \fill[color={rgb,255:red,153;green,153;blue,153}] (1.931,0.479) rectangle ++(0.050,0.050);
  \fill[color={rgb,255:red,153;green,153;blue,153}] (1.960,0.481) rectangle ++(0.050,0.050);
  \fill[color={rgb,255:red,153;green,153;blue,153}] (1.989,0.480) rectangle ++(0.050,0.050);
  \fill[color={rgb,255:red,153;green,153;blue,153}] (2.015,0.471) rectangle ++(0.050,0.050);
  \fill[color={rgb,255:red,153;green,153;blue,153}] (2.047,0.482) rectangle ++(0.050,0.050);
  \fill[color={rgb,255:red,153;green,153;blue,153}] (2.069,0.459) rectangle ++(0.050,0.050);
  \fill[color={rgb,255:red,153;green,153;blue,153}] (2.104,0.485) rectangle ++(0.050,0.050);
  \fill[color={rgb,255:red,153;green,153;blue,153}] (2.128,0.470) rectangle ++(0.050,0.050);
  \fill[color={rgb,255:red,153;green,153;blue,153}] (2.158,0.482) rectangle ++(0.050,0.050);
  \fill[color={rgb,255:red,153;green,153;blue,153}] (2.180,0.456) rectangle ++(0.050,0.050);
  \fill[color={rgb,255:red,153;green,153;blue,153}] (2.211,0.479) rectangle ++(0.050,0.050);
  \fill[color={rgb,255:red,153;green,153;blue,153}] (2.238,0.476) rectangle ++(0.050,0.050);
  \fill[color={rgb,255:red,153;green,153;blue,153}] (2.262,0.460) rectangle ++(0.050,0.050);
  \fill[color={rgb,255:red,153;green,153;blue,153}] (2.289,0.461) rectangle ++(0.050,0.050);
  \fill[color={rgb,255:red,153;green,153;blue,153}] (2.317,0.472) rectangle ++(0.050,0.050);
  \fill[color={rgb,255:red,153;green,153;blue,153}] (2.344,0.473) rectangle ++(0.050,0.050);
  \fill[color={rgb,255:red,153;green,153;blue,153}] (2.369,0.452) rectangle ++(0.050,0.050);
  \fill[color={rgb,255:red,153;green,153;blue,153}] (2.397,0.477) rectangle ++(0.050,0.050);
  \fill[color={rgb,255:red,153;green,153;blue,153}] (2.423,0.485) rectangle ++(0.050,0.050);
  \fill[color={rgb,255:red,153;green,153;blue,153}] (2.449,0.469) rectangle ++(0.050,0.050);
  \fill[color={rgb,255:red,153;green,153;blue,153}] (2.475,0.493) rectangle ++(0.050,0.050);
  \fill[color={rgb,255:red,153;green,153;blue,153}] (2.501,0.492) rectangle ++(0.050,0.050);
  \fill[color={rgb,255:red,153;green,153;blue,153}] (2.526,0.507) rectangle ++(0.050,0.050);
  \fill[color={rgb,255:red,153;green,153;blue,153}] (2.553,0.477) rectangle ++(0.050,0.050);
  \fill[color={rgb,255:red,153;green,153;blue,153}] (2.579,0.488) rectangle ++(0.050,0.050);
  \fill[color={rgb,255:red,153;green,153;blue,153}] (2.607,0.463) rectangle ++(0.050,0.050);
  \fill[color={rgb,255:red,153;green,153;blue,153}] (2.631,0.491) rectangle ++(0.050,0.050);
  \fill[color={rgb,255:red,153;green,153;blue,153}] (2.659,0.474) rectangle ++(0.050,0.050);
  \fill[color={rgb,255:red,153;green,153;blue,153}] (2.686,0.476) rectangle ++(0.050,0.050);
  \fill[color={rgb,255:red,153;green,153;blue,153}] (2.714,0.465) rectangle ++(0.050,0.050);
  \fill[color={rgb,255:red,153;green,153;blue,153}] (2.738,0.486) rectangle ++(0.050,0.050);
  \fill[color={rgb,255:red,153;green,153;blue,153}] (2.767,0.474) rectangle ++(0.050,0.050);
  \fill[color={rgb,255:red,153;green,153;blue,153}] (2.797,0.461) rectangle ++(0.050,0.050);
  \fill[color={rgb,255:red,153;green,153;blue,153}] (2.824,0.465) rectangle ++(0.050,0.050);
  \fill[color={rgb,255:red,153;green,153;blue,153}] (2.855,0.451) rectangle ++(0.050,0.050);
  \fill[color={rgb,255:red,153;green,153;blue,153}] (2.883,0.452) rectangle ++(0.050,0.050);
  \fill[color={rgb,255:red,153;green,153;blue,153}] (2.900,0.492) rectangle ++(0.050,0.050);
  \fill[color={rgb,255:red,153;green,153;blue,153}] (2.938,0.462) rectangle ++(0.050,0.050);
  \fill[color={rgb,255:red,153;green,153;blue,153}] (2.961,0.480) rectangle ++(0.050,0.050);
  \fill[color={rgb,255:red,153;green,153;blue,153}] (2.990,0.479) rectangle ++(0.050,0.050);
  \fill[color={rgb,255:red,153;green,153;blue,153}] (3.019,0.480) rectangle ++(0.050,0.050);
  \fill[color={rgb,255:red,153;green,153;blue,153}] (3.047,0.485) rectangle ++(0.050,0.050);
  \fill[color={rgb,255:red,153;green,153;blue,153}] (3.076,0.486) rectangle ++(0.050,0.050);
  \fill[color={rgb,255:red,153;green,153;blue,153}] (3.113,0.472) rectangle ++(0.050,0.050);
  \fill[color={rgb,255:red,153;green,153;blue,153}] (3.148,0.464) rectangle ++(0.050,0.050);
  \fill[color={rgb,255:red,153;green,153;blue,153}] (3.189,0.443) rectangle ++(0.050,0.050);
  \fill[color={rgb,255:red,153;green,153;blue,153}] (3.201,0.486) rectangle ++(0.050,0.050);
  \fill[color={rgb,255:red,153;green,153;blue,153}] (3.236,0.482) rectangle ++(0.050,0.050);
  \fill[color={rgb,255:red,153;green,153;blue,153}] (3.275,0.471) rectangle ++(0.050,0.050);
  \fill[color={rgb,255:red,153;green,153;blue,153}] (3.310,0.468) rectangle ++(0.050,0.050);
  \fill[color={rgb,255:red,153;green,153;blue,153}] (3.334,0.488) rectangle ++(0.050,0.050);
  \fill[color={rgb,255:red,153;green,153;blue,153}] (3.389,0.454) rectangle ++(0.050,0.050);
  \fill[color={rgb,255:red,153;green,153;blue,153}] (3.443,0.425) rectangle ++(0.050,0.050);
  \fill[color={rgb,255:red,153;green,153;blue,153}] (3.447,0.478) rectangle ++(0.050,0.050);
  \fill[color={rgb,255:red,153;green,153;blue,153}] (3.464,0.509) rectangle ++(0.050,0.050);
  \fill[color={rgb,255:red,153;green,153;blue,153}] (3.541,0.452) rectangle ++(0.050,0.050);
  \fill[color={rgb,255:red,153;green,153;blue,153}] (3.566,0.473) rectangle ++(0.050,0.050);
  \fill[color={rgb,255:red,153;green,153;blue,153}] (3.623,0.452) rectangle ++(0.050,0.050);
  \fill[color={rgb,255:red,153;green,153;blue,153}] (3.647,0.475) rectangle ++(0.050,0.050);
  \fill[color={rgb,255:red,153;green,153;blue,153}] (3.675,0.493) rectangle ++(0.050,0.050);
  \fill[color={rgb,255:red,153;green,153;blue,153}] (3.726,0.484) rectangle ++(0.050,0.050);
  \fill[color={rgb,255:red,153;green,153;blue,153}] (3.770,0.485) rectangle ++(0.050,0.050);
  \fill[color={rgb,255:red,153;green,153;blue,153}] (3.828,0.472) rectangle ++(0.050,0.050);
  \fill[color={rgb,255:red,153;green,153;blue,153}] (3.876,0.472) rectangle ++(0.050,0.050);
  \fill[color={rgb,255:red,153;green,153;blue,153}] (3.916,0.482) rectangle ++(0.050,0.050);
  \fill[color={rgb,255:red,153;green,153;blue,153}] (3.992,0.458) rectangle ++(0.050,0.050);
  \fill[color={rgb,255:red,153;green,153;blue,153}] (4.023,0.480) rectangle ++(0.050,0.050);
  \fill[color={rgb,255:red,153;green,153;blue,153}] (4.075,0.483) rectangle ++(0.050,0.050);
  \fill[color={rgb,255:red,153;green,153;blue,153}] (4.153,0.464) rectangle ++(0.050,0.050);
  \fill[color={rgb,255:red,153;green,153;blue,153}] (4.194,0.481) rectangle ++(0.050,0.050);
  \fill[color={rgb,255:red,153;green,153;blue,153}] (4.257,0.480) rectangle ++(0.050,0.050);
  \fill[color={rgb,255:red,153;green,153;blue,153}] (4.336,0.468) rectangle ++(0.050,0.050);
  \fill[color={rgb,255:red,153;green,153;blue,153}] (4.387,0.481) rectangle ++(0.050,0.050);
  \fill[color={rgb,255:red,153;green,153;blue,153}] (4.465,0.475) rectangle ++(0.050,0.050);
  \fill[color={rgb,255:red,153;green,153;blue,153}] (4.554,0.465) rectangle ++(0.050,0.050);
  \fill[color={rgb,255:red,153;green,153;blue,153}] (4.623,0.471) rectangle ++(0.050,0.050);
  \fill[color={rgb,255:red,153;green,153;blue,153}] (4.695,0.477) rectangle ++(0.050,0.050);
  \fill[color={rgb,255:red,153;green,153;blue,153}] (4.783,0.476) rectangle ++(0.050,0.050);
  \fill[color={rgb,255:red,153;green,153;blue,153}] (4.862,0.483) rectangle ++(0.050,0.050);
  \fill[color={rgb,255:red,153;green,153;blue,153}] (4.961,0.481) rectangle ++(0.050,0.050);
  \fill[color={rgb,255:red,153;green,153;blue,153}] (5.074,0.474) rectangle ++(0.050,0.050);
  \fill[color={rgb,255:red,153;green,153;blue,153}] (5.179,0.476) rectangle ++(0.050,0.050);
  \fill[color={rgb,255:red,153;green,153;blue,153}] (5.308,0.469) rectangle ++(0.050,0.050);
  \fill[color={rgb,255:red,153;green,153;blue,153}] (5.410,0.480) rectangle ++(0.050,0.050);
  \fill[color={rgb,255:red,153;green,153;blue,153}] (5.543,0.478) rectangle ++(0.050,0.050);
  \fill[color={rgb,255:red,153;green,153;blue,153}] (5.690,0.476) rectangle ++(0.050,0.050);
  \fill[color={rgb,255:red,153;green,153;blue,153}] (5.840,0.477) rectangle ++(0.050,0.050);
  \fill[color={rgb,255:red,153;green,153;blue,153}] (5.995,0.481) rectangle ++(0.050,0.050);
  \fill[color={rgb,255:red,153;green,153;blue,153}] (6.196,0.472) rectangle ++(0.050,0.050);
  \fill[color={rgb,255:red,153;green,153;blue,153}] (6.386,0.474) rectangle ++(0.050,0.050);
  \fill[color={rgb,255:red,153;green,153;blue,153}] (6.575,0.483) rectangle ++(0.050,0.050);
  \fill[color={rgb,255:red,153;green,153;blue,153}] (6.832,0.475) rectangle ++(0.050,0.050);
  \fill[color={rgb,255:red,153;green,153;blue,153}] (7.080,0.479) rectangle ++(0.050,0.050);
  \fill[color={rgb,255:red,153;green,153;blue,153}] (7.359,0.482) rectangle ++(0.050,0.050);
  \fill[color={rgb,255:red,153;green,153;blue,153}] (7.713,0.473) rectangle ++(0.050,0.050);
  \fill[color={rgb,255:red,153;green,153;blue,153}] (8.067,0.477) rectangle ++(0.050,0.050);
  \fill[color={rgb,255:red,153;green,153;blue,153}] (8.469,0.481) rectangle ++(0.050,0.050);
  \fill[color={rgb,255:red,153;green,153;blue,153}] (8.984,0.472) rectangle ++(0.050,0.050);
  \fill[color={rgb,255:red,153;green,153;blue,153}] (9.470,0.488) rectangle ++(0.050,0.050);
  \fill[color={rgb,255:red,153;green,153;blue,153}] (9.482,0.613) rectangle ++(0.050,0.050);
  \fill[color={rgb,255:red,153;green,153;blue,153}] (9.490,0.738) rectangle ++(0.050,0.050);
  \fill[color={rgb,255:red,153;green,153;blue,153}] (9.478,0.866) rectangle ++(0.050,0.050);
  \fill[color={rgb,255:red,153;green,153;blue,153}] (9.475,0.991) rectangle ++(0.050,0.050);
  \fill[color={rgb,255:red,153;green,153;blue,153}] (9.470,1.116) rectangle ++(0.050,0.050);
  \fill[color={rgb,255:red,153;green,153;blue,153}] (9.474,1.239) rectangle ++(0.050,0.050);
  \fill[color={rgb,255:red,153;green,153;blue,153}] (9.476,1.363) rectangle ++(0.050,0.050);
  \fill[color={rgb,255:red,153;green,153;blue,153}] (9.471,1.486) rectangle ++(0.050,0.050);
  \fill[color={rgb,255:red,153;green,153;blue,153}] (9.481,1.608) rectangle ++(0.050,0.050);
  \fill[color={rgb,255:red,153;green,153;blue,153}] (9.443,1.732) rectangle ++(0.050,0.050);
  \fill[color={rgb,255:red,153;green,153;blue,153}] (9.480,1.853) rectangle ++(0.050,0.050);
  \fill[color={rgb,255:red,31;green,94;blue,168}] (0.459,0.507) rectangle ++(0.070,0.070);
  \fill[color={rgb,255:red,31;green,94;blue,168}] (0.460,0.454) rectangle ++(0.070,0.070);
  \fill[color={rgb,255:red,31;green,94;blue,168}] (0.527,0.451) rectangle ++(0.070,0.070);
  \fill[color={rgb,255:red,31;green,94;blue,168}] (0.592,0.448) rectangle ++(0.070,0.070);
  \fill[color={rgb,255:red,31;green,94;blue,168}] (0.671,0.460) rectangle ++(0.070,0.070);
  \fill[color={rgb,255:red,31;green,94;blue,168}] (0.758,0.481) rectangle ++(0.070,0.070);
  \fill[color={rgb,255:red,31;green,94;blue,168}] (0.789,0.456) rectangle ++(0.070,0.070);
  \fill[color={rgb,255:red,31;green,94;blue,168}] (0.881,0.488) rectangle ++(0.070,0.070);
  \fill[color={rgb,255:red,31;green,94;blue,168}] (0.928,0.481) rectangle ++(0.070,0.070);
  \fill[color={rgb,255:red,31;green,94;blue,168}] (0.968,0.468) rectangle ++(0.070,0.070);
  \fill[color={rgb,255:red,31;green,94;blue,168}] (1.028,0.477) rectangle ++(0.070,0.070);
  \fill[color={rgb,255:red,31;green,94;blue,168}] (1.065,0.463) rectangle ++(0.070,0.070);
  \fill[color={rgb,255:red,31;green,94;blue,168}] (1.104,0.454) rectangle ++(0.070,0.070);
  \fill[color={rgb,255:red,31;green,94;blue,168}] (1.165,0.470) rectangle ++(0.070,0.070);
  \fill[color={rgb,255:red,31;green,94;blue,168}] (1.218,0.479) rectangle ++(0.070,0.070);
  \fill[color={rgb,255:red,31;green,94;blue,168}] (1.250,0.464) rectangle ++(0.070,0.070);
  \fill[color={rgb,255:red,31;green,94;blue,168}] (1.284,0.454) rectangle ++(0.070,0.070);
  \fill[color={rgb,255:red,31;green,94;blue,168}] (1.351,0.486) rectangle ++(0.070,0.070);
  \fill[color={rgb,255:red,31;green,94;blue,168}] (1.381,0.473) rectangle ++(0.070,0.070);
  \fill[color={rgb,255:red,31;green,94;blue,168}] (1.418,0.470) rectangle ++(0.070,0.070);
  \fill[color={rgb,255:red,31;green,94;blue,168}] (1.473,0.495) rectangle ++(0.070,0.070);
  \fill[color={rgb,255:red,31;green,94;blue,168}] (1.490,0.464) rectangle ++(0.070,0.070);
  \fill[color={rgb,255:red,31;green,94;blue,168}] (1.543,0.489) rectangle ++(0.070,0.070);
  \fill[color={rgb,255:red,31;green,94;blue,168}] (1.564,0.466) rectangle ++(0.070,0.070);
  \fill[color={rgb,255:red,31;green,94;blue,168}] (1.615,0.493) rectangle ++(0.070,0.070);
  \fill[color={rgb,255:red,31;green,94;blue,168}] (1.636,0.469) rectangle ++(0.070,0.070);
  \fill[color={rgb,255:red,31;green,94;blue,168}] (1.654,0.439) rectangle ++(0.070,0.070);
  \fill[color={rgb,255:red,31;green,94;blue,168}] (1.707,0.478) rectangle ++(0.070,0.070);
  \fill[color={rgb,255:red,31;green,94;blue,168}] (1.725,0.447) rectangle ++(0.070,0.070);
  \fill[color={rgb,255:red,31;green,94;blue,168}] (1.770,0.475) rectangle ++(0.070,0.070);
  \fill[color={rgb,255:red,31;green,94;blue,168}] (1.798,0.467) rectangle ++(0.070,0.070);
  \fill[color={rgb,255:red,31;green,94;blue,168}] (1.837,0.486) rectangle ++(0.070,0.070);
  \fill[color={rgb,255:red,31;green,94;blue,168}] (1.875,0.504) rectangle ++(0.070,0.070);
  \fill[color={rgb,255:red,31;green,94;blue,168}] (1.887,0.459) rectangle ++(0.070,0.070);
  \fill[color={rgb,255:red,31;green,94;blue,168}] (1.921,0.469) rectangle ++(0.070,0.070);
  \fill[color={rgb,255:red,31;green,94;blue,168}] (1.950,0.471) rectangle ++(0.070,0.070);
  \fill[color={rgb,255:red,31;green,94;blue,168}] (1.979,0.470) rectangle ++(0.070,0.070);
  \fill[color={rgb,255:red,31;green,94;blue,168}] (2.005,0.461) rectangle ++(0.070,0.070);
  \fill[color={rgb,255:red,31;green,94;blue,168}] (2.037,0.472) rectangle ++(0.070,0.070);
  \fill[color={rgb,255:red,31;green,94;blue,168}] (2.059,0.449) rectangle ++(0.070,0.070);
  \fill[color={rgb,255:red,31;green,94;blue,168}] (2.094,0.475) rectangle ++(0.070,0.070);
  \fill[color={rgb,255:red,31;green,94;blue,168}] (2.118,0.460) rectangle ++(0.070,0.070);
  \fill[color={rgb,255:red,31;green,94;blue,168}] (2.148,0.472) rectangle ++(0.070,0.070);
  \fill[color={rgb,255:red,31;green,94;blue,168}] (2.170,0.446) rectangle ++(0.070,0.070);
  \fill[color={rgb,255:red,31;green,94;blue,168}] (2.201,0.469) rectangle ++(0.070,0.070);
  \fill[color={rgb,255:red,31;green,94;blue,168}] (2.228,0.466) rectangle ++(0.070,0.070);
  \fill[color={rgb,255:red,31;green,94;blue,168}] (2.252,0.450) rectangle ++(0.070,0.070);
  \fill[color={rgb,255:red,31;green,94;blue,168}] (2.279,0.451) rectangle ++(0.070,0.070);
  \fill[color={rgb,255:red,31;green,94;blue,168}] (2.307,0.462) rectangle ++(0.070,0.070);
  \fill[color={rgb,255:red,31;green,94;blue,168}] (2.334,0.463) rectangle ++(0.070,0.070);
  \fill[color={rgb,255:red,31;green,94;blue,168}] (2.359,0.442) rectangle ++(0.070,0.070);
  \fill[color={rgb,255:red,31;green,94;blue,168}] (2.387,0.467) rectangle ++(0.070,0.070);
  \fill[color={rgb,255:red,31;green,94;blue,168}] (2.413,0.475) rectangle ++(0.070,0.070);
  \fill[color={rgb,255:red,31;green,94;blue,168}] (2.439,0.459) rectangle ++(0.070,0.070);
  \fill[color={rgb,255:red,31;green,94;blue,168}] (2.465,0.483) rectangle ++(0.070,0.070);
  \fill[color={rgb,255:red,31;green,94;blue,168}] (2.491,0.482) rectangle ++(0.070,0.070);
  \fill[color={rgb,255:red,31;green,94;blue,168}] (2.516,0.497) rectangle ++(0.070,0.070);
  \fill[color={rgb,255:red,31;green,94;blue,168}] (2.543,0.467) rectangle ++(0.070,0.070);
  \fill[color={rgb,255:red,31;green,94;blue,168}] (2.569,0.478) rectangle ++(0.070,0.070);
  \fill[color={rgb,255:red,31;green,94;blue,168}] (2.597,0.453) rectangle ++(0.070,0.070);
  \fill[color={rgb,255:red,31;green,94;blue,168}] (2.621,0.481) rectangle ++(0.070,0.070);
  \fill[color={rgb,255:red,31;green,94;blue,168}] (2.649,0.464) rectangle ++(0.070,0.070);
  \fill[color={rgb,255:red,31;green,94;blue,168}] (2.676,0.466) rectangle ++(0.070,0.070);
  \fill[color={rgb,255:red,31;green,94;blue,168}] (2.704,0.455) rectangle ++(0.070,0.070);
  \fill[color={rgb,255:red,31;green,94;blue,168}] (2.728,0.476) rectangle ++(0.070,0.070);
  \fill[color={rgb,255:red,31;green,94;blue,168}] (2.757,0.464) rectangle ++(0.070,0.070);
  \fill[color={rgb,255:red,31;green,94;blue,168}] (2.787,0.451) rectangle ++(0.070,0.070);
  \fill[color={rgb,255:red,31;green,94;blue,168}] (2.814,0.455) rectangle ++(0.070,0.070);
  \fill[color={rgb,255:red,31;green,94;blue,168}] (2.845,0.441) rectangle ++(0.070,0.070);
  \fill[color={rgb,255:red,31;green,94;blue,168}] (2.873,0.442) rectangle ++(0.070,0.070);
  \fill[color={rgb,255:red,31;green,94;blue,168}] (2.890,0.482) rectangle ++(0.070,0.070);
  \fill[color={rgb,255:red,31;green,94;blue,168}] (2.928,0.452) rectangle ++(0.070,0.070);
  \fill[color={rgb,255:red,31;green,94;blue,168}] (2.951,0.470) rectangle ++(0.070,0.070);
  \fill[color={rgb,255:red,31;green,94;blue,168}] (2.980,0.469) rectangle ++(0.070,0.070);
  \fill[color={rgb,255:red,31;green,94;blue,168}] (3.009,0.470) rectangle ++(0.070,0.070);
  \fill[color={rgb,255:red,31;green,94;blue,168}] (3.037,0.475) rectangle ++(0.070,0.070);
  \fill[color={rgb,255:red,31;green,94;blue,168}] (3.066,0.476) rectangle ++(0.070,0.070);
  \fill[color={rgb,255:red,31;green,94;blue,168}] (3.103,0.462) rectangle ++(0.070,0.070);
  \fill[color={rgb,255:red,31;green,94;blue,168}] (3.138,0.454) rectangle ++(0.070,0.070);
  \fill[color={rgb,255:red,31;green,94;blue,168}] (3.179,0.433) rectangle ++(0.070,0.070);
  \fill[color={rgb,255:red,31;green,94;blue,168}] (3.191,0.476) rectangle ++(0.070,0.070);
  \fill[color={rgb,255:red,31;green,94;blue,168}] (3.226,0.472) rectangle ++(0.070,0.070);
  \fill[color={rgb,255:red,31;green,94;blue,168}] (3.265,0.461) rectangle ++(0.070,0.070);
  \fill[color={rgb,255:red,31;green,94;blue,168}] (3.300,0.458) rectangle ++(0.070,0.070);
  \fill[color={rgb,255:red,31;green,94;blue,168}] (3.324,0.478) rectangle ++(0.070,0.070);
  \fill[color={rgb,255:red,31;green,94;blue,168}] (3.379,0.444) rectangle ++(0.070,0.070);
  \fill[color={rgb,255:red,31;green,94;blue,168}] (3.437,0.468) rectangle ++(0.070,0.070);
  \fill[color={rgb,255:red,31;green,94;blue,168}] (3.454,0.499) rectangle ++(0.070,0.070);
  \fill[color={rgb,255:red,31;green,94;blue,168}] (3.531,0.442) rectangle ++(0.070,0.070);
  \fill[color={rgb,255:red,31;green,94;blue,168}] (3.556,0.463) rectangle ++(0.070,0.070);
  \fill[color={rgb,255:red,31;green,94;blue,168}] (3.613,0.442) rectangle ++(0.070,0.070);
  \fill[color={rgb,255:red,31;green,94;blue,168}] (3.637,0.465) rectangle ++(0.070,0.070);
  \fill[color={rgb,255:red,31;green,94;blue,168}] (3.665,0.483) rectangle ++(0.070,0.070);
  \fill[color={rgb,255:red,31;green,94;blue,168}] (3.716,0.474) rectangle ++(0.070,0.070);
  \fill[color={rgb,255:red,31;green,94;blue,168}] (3.760,0.475) rectangle ++(0.070,0.070);
  \fill[color={rgb,255:red,31;green,94;blue,168}] (3.818,0.462) rectangle ++(0.070,0.070);
  \fill[color={rgb,255:red,31;green,94;blue,168}] (3.866,0.462) rectangle ++(0.070,0.070);
  \fill[color={rgb,255:red,31;green,94;blue,168}] (3.906,0.472) rectangle ++(0.070,0.070);
  \fill[color={rgb,255:red,31;green,94;blue,168}] (3.982,0.448) rectangle ++(0.070,0.070);
  \fill[color={rgb,255:red,31;green,94;blue,168}] (4.013,0.470) rectangle ++(0.070,0.070);
  \fill[color={rgb,255:red,31;green,94;blue,168}] (4.065,0.473) rectangle ++(0.070,0.070);
  \fill[color={rgb,255:red,31;green,94;blue,168}] (4.143,0.454) rectangle ++(0.070,0.070);
  \fill[color={rgb,255:red,31;green,94;blue,168}] (4.184,0.471) rectangle ++(0.070,0.070);
  \fill[color={rgb,255:red,31;green,94;blue,168}] (4.247,0.470) rectangle ++(0.070,0.070);
  \fill[color={rgb,255:red,31;green,94;blue,168}] (4.326,0.458) rectangle ++(0.070,0.070);
  \fill[color={rgb,255:red,31;green,94;blue,168}] (4.377,0.471) rectangle ++(0.070,0.070);
  \fill[color={rgb,255:red,31;green,94;blue,168}] (4.455,0.465) rectangle ++(0.070,0.070);
  \fill[color={rgb,255:red,31;green,94;blue,168}] (4.544,0.455) rectangle ++(0.070,0.070);
  \fill[color={rgb,255:red,31;green,94;blue,168}] (4.613,0.461) rectangle ++(0.070,0.070);
  \fill[color={rgb,255:red,31;green,94;blue,168}] (4.685,0.467) rectangle ++(0.070,0.070);
  \fill[color={rgb,255:red,31;green,94;blue,168}] (4.773,0.466) rectangle ++(0.070,0.070);
  \fill[color={rgb,255:red,31;green,94;blue,168}] (4.852,0.473) rectangle ++(0.070,0.070);
  \fill[color={rgb,255:red,31;green,94;blue,168}] (4.951,0.471) rectangle ++(0.070,0.070);
  \fill[color={rgb,255:red,31;green,94;blue,168}] (5.064,0.464) rectangle ++(0.070,0.070);
  \fill[color={rgb,255:red,31;green,94;blue,168}] (5.169,0.466) rectangle ++(0.070,0.070);
  \fill[color={rgb,255:red,31;green,94;blue,168}] (5.298,0.459) rectangle ++(0.070,0.070);
  \fill[color={rgb,255:red,31;green,94;blue,168}] (5.400,0.470) rectangle ++(0.070,0.070);
  \fill[color={rgb,255:red,31;green,94;blue,168}] (5.533,0.468) rectangle ++(0.070,0.070);
  \fill[color={rgb,255:red,31;green,94;blue,168}] (5.680,0.466) rectangle ++(0.070,0.070);
  \fill[color={rgb,255:red,31;green,94;blue,168}] (5.830,0.467) rectangle ++(0.070,0.070);
  \fill[color={rgb,255:red,31;green,94;blue,168}] (5.985,0.471) rectangle ++(0.070,0.070);
  \fill[color={rgb,255:red,31;green,94;blue,168}] (6.186,0.462) rectangle ++(0.070,0.070);
  \fill[color={rgb,255:red,31;green,94;blue,168}] (6.376,0.464) rectangle ++(0.070,0.070);
  \fill[color={rgb,255:red,31;green,94;blue,168}] (6.565,0.473) rectangle ++(0.070,0.070);
  \fill[color={rgb,255:red,31;green,94;blue,168}] (6.822,0.465) rectangle ++(0.070,0.070);
  \fill[color={rgb,255:red,31;green,94;blue,168}] (7.070,0.469) rectangle ++(0.070,0.070);
  \fill[color={rgb,255:red,31;green,94;blue,168}] (7.349,0.472) rectangle ++(0.070,0.070);
  \fill[color={rgb,255:red,31;green,94;blue,168}] (7.703,0.463) rectangle ++(0.070,0.070);
  \fill[color={rgb,255:red,31;green,94;blue,168}] (8.057,0.467) rectangle ++(0.070,0.070);
  \fill[color={rgb,255:red,31;green,94;blue,168}] (8.459,0.471) rectangle ++(0.070,0.070);
  \fill[color={rgb,255:red,31;green,94;blue,168}] (8.974,0.462) rectangle ++(0.070,0.070);
  \fill[color={rgb,255:red,31;green,94;blue,168}] (9.460,0.478) rectangle ++(0.070,0.070);
  \draw[color={rgb,255:red,31;green,94;blue,168}, line width=0.53pt] (0.494,0.503) -- (9.495,0.500);
  \fill[color={rgb,255:red,51;green,161;blue,74}] (0.825,6.472) rectangle ++(0.070,0.070);
  \fill[color={rgb,255:red,51;green,161;blue,74}] (0.479,3.408) rectangle ++(0.070,0.070);
  \fill[color={rgb,255:red,51;green,161;blue,74}] (0.480,3.304) rectangle ++(0.070,0.070);
  \fill[color={rgb,255:red,51;green,161;blue,74}] (0.497,3.243) rectangle ++(0.070,0.070);
  \fill[color={rgb,255:red,51;green,161;blue,74}] (0.461,3.217) rectangle ++(0.070,0.070);
  \fill[color={rgb,255:red,51;green,161;blue,74}] (0.479,3.112) rectangle ++(0.070,0.070);
  \fill[color={rgb,255:red,51;green,161;blue,74}] (0.459,3.077) rectangle ++(0.070,0.070);
  \fill[color={rgb,255:red,51;green,161;blue,74}] (0.462,3.030) rectangle ++(0.070,0.070);
  \fill[color={rgb,255:red,51;green,161;blue,74}] (0.444,2.995) rectangle ++(0.070,0.070);
  \fill[color={rgb,255:red,51;green,161;blue,74}] (0.481,2.932) rectangle ++(0.070,0.070);
  \fill[color={rgb,255:red,51;green,161;blue,74}] (0.461,2.900) rectangle ++(0.070,0.070);
  \fill[color={rgb,255:red,51;green,161;blue,74}] (0.461,2.857) rectangle ++(0.070,0.070);
  \fill[color={rgb,255:red,51;green,161;blue,74}] (0.438,2.826) rectangle ++(0.070,0.070);
  \fill[color={rgb,255:red,51;green,161;blue,74}] (0.449,2.779) rectangle ++(0.070,0.070);
  \fill[color={rgb,255:red,51;green,161;blue,74}] (0.458,2.735) rectangle ++(0.070,0.070);
  \fill[color={rgb,255:red,51;green,161;blue,74}] (0.457,2.696) rectangle ++(0.070,0.070);
  \fill[color={rgb,255:red,51;green,161;blue,74}] (0.461,2.655) rectangle ++(0.070,0.070);
  \fill[color={rgb,255:red,51;green,161;blue,74}] (0.455,2.618) rectangle ++(0.070,0.070);
  \fill[color={rgb,255:red,51;green,161;blue,74}] (0.452,2.580) rectangle ++(0.070,0.070);
  \fill[color={rgb,255:red,51;green,161;blue,74}] (0.448,2.505) rectangle ++(0.070,0.070);
  \fill[color={rgb,255:red,51;green,161;blue,74}] (0.465,2.464) rectangle ++(0.070,0.070);
  \fill[color={rgb,255:red,51;green,161;blue,74}] (0.463,2.427) rectangle ++(0.070,0.070);
  \fill[color={rgb,255:red,51;green,161;blue,74}] (0.457,2.355) rectangle ++(0.070,0.070);
  \fill[color={rgb,255:red,51;green,161;blue,74}] (0.447,2.321) rectangle ++(0.070,0.070);
  \fill[color={rgb,255:red,51;green,161;blue,74}] (0.455,2.248) rectangle ++(0.070,0.070);
  \fill[color={rgb,255:red,51;green,161;blue,74}] (0.434,2.178) rectangle ++(0.070,0.070);
  \fill[color={rgb,255:red,51;green,161;blue,74}] (0.442,2.071) rectangle ++(0.070,0.070);
  \draw[color={rgb,255:red,51;green,161;blue,74}, line width=0.53pt] (0.434,2.110) -- (0.838,6.509);
  \fill[color={rgb,255:red,242;green,140;blue,26}] (7.828,6.465) rectangle ++(0.070,0.070);
  \fill[color={rgb,255:red,242;green,140;blue,26}] (7.645,6.468) rectangle ++(0.070,0.070);
  \fill[color={rgb,255:red,242;green,140;blue,26}] (7.480,6.480) rectangle ++(0.070,0.070);
  \fill[color={rgb,255:red,242;green,140;blue,26}] (7.272,6.448) rectangle ++(0.070,0.070);
  \fill[color={rgb,255:red,242;green,140;blue,26}] (7.116,6.457) rectangle ++(0.070,0.070);
  \fill[color={rgb,255:red,242;green,140;blue,26}] (6.959,6.459) rectangle ++(0.070,0.070);
  \fill[color={rgb,255:red,242;green,140;blue,26}] (6.815,6.470) rectangle ++(0.070,0.070);
  \fill[color={rgb,255:red,242;green,140;blue,26}] (6.670,6.474) rectangle ++(0.070,0.070);
  \fill[color={rgb,255:red,242;green,140;blue,26}] (6.533,6.483) rectangle ++(0.070,0.070);
  \fill[color={rgb,255:red,242;green,140;blue,26}] (6.385,6.474) rectangle ++(0.070,0.070);
  \fill[color={rgb,255:red,242;green,140;blue,26}] (6.243,6.467) rectangle ++(0.070,0.070);
  \fill[color={rgb,255:red,242;green,140;blue,26}] (6.115,6.473) rectangle ++(0.070,0.070);
  \fill[color={rgb,255:red,242;green,140;blue,26}] (5.977,6.460) rectangle ++(0.070,0.070);
  \fill[color={rgb,255:red,242;green,140;blue,26}] (5.860,6.470) rectangle ++(0.070,0.070);
  \fill[color={rgb,255:red,242;green,140;blue,26}] (5.733,6.463) rectangle ++(0.070,0.070);
  \fill[color={rgb,255:red,242;green,140;blue,26}] (5.608,6.454) rectangle ++(0.070,0.070);
  \fill[color={rgb,255:red,242;green,140;blue,26}] (5.499,6.463) rectangle ++(0.070,0.070);
  \fill[color={rgb,255:red,242;green,140;blue,26}] (5.376,6.448) rectangle ++(0.070,0.070);
  \fill[color={rgb,255:red,242;green,140;blue,26}] (5.271,6.455) rectangle ++(0.070,0.070);
  \fill[color={rgb,255:red,242;green,140;blue,26}] (5.158,6.446) rectangle ++(0.070,0.070);
  \fill[color={rgb,255:red,242;green,140;blue,26}] (5.075,6.486) rectangle ++(0.070,0.070);
  \fill[color={rgb,255:red,242;green,140;blue,26}] (4.963,6.471) rectangle ++(0.070,0.070);
  \fill[color={rgb,255:red,242;green,140;blue,26}] (4.853,6.457) rectangle ++(0.070,0.070);
  \fill[color={rgb,255:red,242;green,140;blue,26}] (4.768,6.484) rectangle ++(0.070,0.070);
  \fill[color={rgb,255:red,242;green,140;blue,26}] (4.657,6.460) rectangle ++(0.070,0.070);
  \fill[color={rgb,255:red,242;green,140;blue,26}] (4.557,6.451) rectangle ++(0.070,0.070);
  \fill[color={rgb,255:red,242;green,140;blue,26}] (4.466,6.459) rectangle ++(0.070,0.070);
  \fill[color={rgb,255:red,242;green,140;blue,26}] (4.373,6.460) rectangle ++(0.070,0.070);
  \fill[color={rgb,255:red,242;green,140;blue,26}] (4.274,6.442) rectangle ++(0.070,0.070);
  \fill[color={rgb,255:red,242;green,140;blue,26}] (4.192,6.465) rectangle ++(0.070,0.070);
  \fill[color={rgb,255:red,242;green,140;blue,26}] (4.095,6.444) rectangle ++(0.070,0.070);
  \fill[color={rgb,255:red,242;green,140;blue,26}] (4.018,6.476) rectangle ++(0.070,0.070);
  \fill[color={rgb,255:red,242;green,140;blue,26}] (3.927,6.463) rectangle ++(0.070,0.070);
  \fill[color={rgb,255:red,242;green,140;blue,26}] (3.844,6.475) rectangle ++(0.070,0.070);
  \fill[color={rgb,255:red,242;green,140;blue,26}] (3.760,6.480) rectangle ++(0.070,0.070);
  \fill[color={rgb,255:red,242;green,140;blue,26}] (3.670,6.462) rectangle ++(0.070,0.070);
  \fill[color={rgb,255:red,242;green,140;blue,26}] (3.585,6.457) rectangle ++(0.070,0.070);
  \fill[color={rgb,255:red,242;green,140;blue,26}] (3.497,6.435) rectangle ++(0.070,0.070);
  \fill[color={rgb,255:red,242;green,140;blue,26}] (3.419,6.452) rectangle ++(0.070,0.070);
  \fill[color={rgb,255:red,242;green,140;blue,26}] (3.342,6.478) rectangle ++(0.070,0.070);
  \fill[color={rgb,255:red,242;green,140;blue,26}] (3.259,6.466) rectangle ++(0.070,0.070);
  \fill[color={rgb,255:red,242;green,140;blue,26}] (3.177,6.462) rectangle ++(0.070,0.070);
  \fill[color={rgb,255:red,242;green,140;blue,26}] (3.099,6.476) rectangle ++(0.070,0.070);
  \fill[color={rgb,255:red,242;green,140;blue,26}] (3.020,6.486) rectangle ++(0.070,0.070);
  \fill[color={rgb,255:red,242;green,140;blue,26}] (2.936,6.447) rectangle ++(0.070,0.070);
  \fill[color={rgb,255:red,242;green,140;blue,26}] (2.859,6.474) rectangle ++(0.070,0.070);
  \fill[color={rgb,255:red,242;green,140;blue,26}] (2.779,6.461) rectangle ++(0.070,0.070);
  \fill[color={rgb,255:red,242;green,140;blue,26}] (2.701,6.468) rectangle ++(0.070,0.070);
  \fill[color={rgb,255:red,242;green,140;blue,26}] (2.623,6.496) rectangle ++(0.070,0.070);
  \fill[color={rgb,255:red,242;green,140;blue,26}] (2.544,6.468) rectangle ++(0.070,0.070);
  \fill[color={rgb,255:red,242;green,140;blue,26}] (2.465,6.450) rectangle ++(0.070,0.070);
  \fill[color={rgb,255:red,242;green,140;blue,26}] (2.387,6.444) rectangle ++(0.070,0.070);
  \fill[color={rgb,255:red,242;green,140;blue,26}] (2.308,6.452) rectangle ++(0.070,0.070);
  \fill[color={rgb,255:red,242;green,140;blue,26}] (2.228,6.491) rectangle ++(0.070,0.070);
  \fill[color={rgb,255:red,242;green,140;blue,26}] (2.153,6.432) rectangle ++(0.070,0.070);
  \fill[color={rgb,255:red,242;green,140;blue,26}] (2.072,6.457) rectangle ++(0.070,0.070);
  \fill[color={rgb,255:red,242;green,140;blue,26}] (1.990,6.487) rectangle ++(0.070,0.070);
  \fill[color={rgb,255:red,242;green,140;blue,26}] (1.912,6.468) rectangle ++(0.070,0.070);
  \fill[color={rgb,255:red,242;green,140;blue,26}] (1.834,6.458) rectangle ++(0.070,0.070);
  \fill[color={rgb,255:red,242;green,140;blue,26}] (1.756,6.440) rectangle ++(0.070,0.070);
  \fill[color={rgb,255:red,242;green,140;blue,26}] (1.671,6.468) rectangle ++(0.070,0.070);
  \fill[color={rgb,255:red,242;green,140;blue,26}] (1.592,6.456) rectangle ++(0.070,0.070);
  \fill[color={rgb,255:red,242;green,140;blue,26}] (1.510,6.457) rectangle ++(0.070,0.070);
  \fill[color={rgb,255:red,242;green,140;blue,26}] (1.420,6.492) rectangle ++(0.070,0.070);
  \fill[color={rgb,255:red,242;green,140;blue,26}] (1.345,6.458) rectangle ++(0.070,0.070);
  \fill[color={rgb,255:red,242;green,140;blue,26}] (1.262,6.456) rectangle ++(0.070,0.070);
  \fill[color={rgb,255:red,242;green,140;blue,26}] (1.165,6.497) rectangle ++(0.070,0.070);
  \fill[color={rgb,255:red,242;green,140;blue,26}] (1.089,6.467) rectangle ++(0.070,0.070);
  \fill[color={rgb,255:red,242;green,140;blue,26}] (1.005,6.458) rectangle ++(0.070,0.070);
  \fill[color={rgb,255:red,242;green,140;blue,26}] (0.917,6.461) rectangle ++(0.070,0.070);
  \fill[color={rgb,255:red,242;green,140;blue,26}] (0.729,6.486) rectangle ++(0.070,0.070);
  \fill[color={rgb,255:red,242;green,140;blue,26}] (0.645,6.471) rectangle ++(0.070,0.070);
  \fill[color={rgb,255:red,242;green,140;blue,26}] (0.556,6.462) rectangle ++(0.070,0.070);
  \fill[color={rgb,255:red,242;green,140;blue,26}] (0.463,6.460) rectangle ++(0.070,0.070);
  \draw[color={rgb,255:red,242;green,140;blue,26}, line width=0.53pt] (0.498,6.500) -- (7.863,6.499);
  \draw[color={rgb,255:red,140;green,76;blue,178}, line width=0.53pt] (-1.106,-3.248) -- (1.987,8.347);
  \draw[color={rgb,255:red,199;green,15;blue,46}, line width=0.88pt, line join=round] (-0.105,0.503) -- (0.286,0.503) -- (1.494,6.500) -- (0.837,6.500) -- (-0.105,0.503);
  \fill[color={rgb,255:red,199;green,15;blue,46}] (-0.105,0.503) circle (0.075);
  \fill[color={rgb,255:red,199;green,15;blue,46}] (0.286,0.503) circle (0.075);
  \fill[color={rgb,255:red,199;green,15;blue,46}] (1.494,6.500) circle (0.075);
  \fill[color={rgb,255:red,199;green,15;blue,46}] (0.837,6.500) circle (0.075);
  \node[color={rgb,255:red,255;green,255;blue,255}, inner sep=0pt] at (0.090,0.503) {0.39 m};
  \node[color={rgb,255:red,255;green,255;blue,255}, inner sep=0pt] at (0.890,3.501) {6.12 m};
  \node[color={rgb,255:red,255;green,255;blue,255}, inner sep=0pt] at (1.165,6.500) {0.66 m};
  \node[color={rgb,255:red,255;green,255;blue,255}, inner sep=0pt] at (0.366,3.501) {6.07 m};
\end{tikzpicture}
}\\[2pt]
    {\small\thema{-p hesse:=false}, \(y = m x + t\)}
  \end{minipage}\hspace{0.05\linewidth}
  \begin{minipage}[t]{0.4\linewidth}\centering
    \resizebox{\linewidth}{!}{% made by tools/make_rviz_figures.py -- do not edit
\begin{tikzpicture}[x=0.62cm, y=0.62cm, font=\scriptsize]
  \fill[color={rgb,255:red,33;green,36;blue,43}] (-0.4,-0.4) rectangle (10.4,7.4);
  \clip (-0.4,-0.4) rectangle (10.4,7.4);
  \draw[color={rgb,255:red,68;green,73;blue,84}, line width=0.53pt, shift={(0,0.5)}] (-1,-3) grid (11,9);
  \draw[color={rgb,255:red,230;green,40;blue,40}, line width=1.2pt] (0.000,0.000) -- (0.600,0.000);
  \draw[color={rgb,255:red,40;green,200;blue,60}, line width=1.2pt] (0.000,0.000) -- (0.000,0.600);
  \draw[color={rgb,255:red,230;green,40;blue,40}, line width=1.2pt] (2.500,2.000) -- (3.100,2.000);
  \draw[color={rgb,255:red,40;green,200;blue,60}, line width=1.2pt] (2.500,2.000) -- (2.500,2.600);
  \fill[color={rgb,255:red,153;green,153;blue,153}] (9.484,1.975) rectangle ++(0.050,0.050);
  \fill[color={rgb,255:red,153;green,153;blue,153}] (9.462,2.097) rectangle ++(0.050,0.050);
  \fill[color={rgb,255:red,153;green,153;blue,153}] (9.465,2.219) rectangle ++(0.050,0.050);
  \fill[color={rgb,255:red,153;green,153;blue,153}] (9.489,2.343) rectangle ++(0.050,0.050);
  \fill[color={rgb,255:red,153;green,153;blue,153}] (9.486,2.465) rectangle ++(0.050,0.050);
  \fill[color={rgb,255:red,153;green,153;blue,153}] (9.482,2.588) rectangle ++(0.050,0.050);
  \fill[color={rgb,255:red,153;green,153;blue,153}] (9.477,2.711) rectangle ++(0.050,0.050);
  \fill[color={rgb,255:red,153;green,153;blue,153}] (9.466,2.833) rectangle ++(0.050,0.050);
  \fill[color={rgb,255:red,153;green,153;blue,153}] (9.515,2.964) rectangle ++(0.050,0.050);
  \fill[color={rgb,255:red,153;green,153;blue,153}] (9.494,3.087) rectangle ++(0.050,0.050);
  \fill[color={rgb,255:red,153;green,153;blue,153}] (9.465,3.208) rectangle ++(0.050,0.050);
  \fill[color={rgb,255:red,153;green,153;blue,153}] (9.482,3.337) rectangle ++(0.050,0.050);
  \fill[color={rgb,255:red,153;green,153;blue,153}] (9.500,3.468) rectangle ++(0.050,0.050);
  \fill[color={rgb,255:red,153;green,153;blue,153}] (9.490,3.595) rectangle ++(0.050,0.050);
  \fill[color={rgb,255:red,153;green,153;blue,153}] (9.481,3.722) rectangle ++(0.050,0.050);
  \fill[color={rgb,255:red,153;green,153;blue,153}] (9.511,3.860) rectangle ++(0.050,0.050);
  \fill[color={rgb,255:red,153;green,153;blue,153}] (9.482,3.984) rectangle ++(0.050,0.050);
  \fill[color={rgb,255:red,153;green,153;blue,153}] (9.469,4.113) rectangle ++(0.050,0.050);
  \fill[color={rgb,255:red,153;green,153;blue,153}] (9.478,4.251) rectangle ++(0.050,0.050);
  \fill[color={rgb,255:red,153;green,153;blue,153}] (9.460,4.380) rectangle ++(0.050,0.050);
  \fill[color={rgb,255:red,153;green,153;blue,153}] (6.661,3.499) rectangle ++(0.050,0.050);
  \fill[color={rgb,255:red,153;green,153;blue,153}] (6.409,3.485) rectangle ++(0.050,0.050);
  \fill[color={rgb,255:red,153;green,153;blue,153}] (6.269,3.508) rectangle ++(0.050,0.050);
  \fill[color={rgb,255:red,153;green,153;blue,153}] (6.162,3.540) rectangle ++(0.050,0.050);
  \fill[color={rgb,255:red,153;green,153;blue,153}] (6.106,3.592) rectangle ++(0.050,0.050);
  \fill[color={rgb,255:red,153;green,153;blue,153}] (6.017,3.627) rectangle ++(0.050,0.050);
  \fill[color={rgb,255:red,153;green,153;blue,153}] (5.941,3.666) rectangle ++(0.050,0.050);
  \fill[color={rgb,255:red,153;green,153;blue,153}] (5.907,3.724) rectangle ++(0.050,0.050);
  \fill[color={rgb,255:red,153;green,153;blue,153}] (5.857,3.773) rectangle ++(0.050,0.050);
  \fill[color={rgb,255:red,153;green,153;blue,153}] (5.816,3.827) rectangle ++(0.050,0.050);
  \fill[color={rgb,255:red,153;green,153;blue,153}] (5.799,3.894) rectangle ++(0.050,0.050);
  \fill[color={rgb,255:red,153;green,153;blue,153}] (5.755,3.946) rectangle ++(0.050,0.050);
  \fill[color={rgb,255:red,153;green,153;blue,153}] (5.742,4.016) rectangle ++(0.050,0.050);
  \fill[color={rgb,255:red,153;green,153;blue,153}] (5.709,4.075) rectangle ++(0.050,0.050);
  \fill[color={rgb,255:red,153;green,153;blue,153}] (5.661,4.124) rectangle ++(0.050,0.050);
  \fill[color={rgb,255:red,153;green,153;blue,153}] (5.669,4.212) rectangle ++(0.050,0.050);
  \fill[color={rgb,255:red,153;green,153;blue,153}] (5.674,4.299) rectangle ++(0.050,0.050);
  \fill[color={rgb,255:red,153;green,153;blue,153}] (5.684,4.393) rectangle ++(0.050,0.050);
  \fill[color={rgb,255:red,153;green,153;blue,153}] (5.716,4.507) rectangle ++(0.050,0.050);
  \fill[color={rgb,255:red,153;green,153;blue,153}] (5.753,4.629) rectangle ++(0.050,0.050);
  \fill[color={rgb,255:red,153;green,153;blue,153}] (7.825,6.465) rectangle ++(0.050,0.050);
  \fill[color={rgb,255:red,153;green,153;blue,153}] (7.662,6.484) rectangle ++(0.050,0.050);
  \fill[color={rgb,255:red,153;green,153;blue,153}] (7.481,6.482) rectangle ++(0.050,0.050);
  \fill[color={rgb,255:red,153;green,153;blue,153}] (7.296,6.471) rectangle ++(0.050,0.050);
  \fill[color={rgb,255:red,153;green,153;blue,153}] (7.113,6.454) rectangle ++(0.050,0.050);
  \fill[color={rgb,255:red,153;green,153;blue,153}] (6.949,6.449) rectangle ++(0.050,0.050);
  \fill[color={rgb,255:red,153;green,153;blue,153}] (6.828,6.482) rectangle ++(0.050,0.050);
  \fill[color={rgb,255:red,153;green,153;blue,153}] (6.688,6.493) rectangle ++(0.050,0.050);
  \fill[color={rgb,255:red,153;green,153;blue,153}] (6.510,6.457) rectangle ++(0.050,0.050);
  \fill[color={rgb,255:red,153;green,153;blue,153}] (6.386,6.474) rectangle ++(0.050,0.050);
  \fill[color={rgb,255:red,153;green,153;blue,153}] (6.255,6.480) rectangle ++(0.050,0.050);
  \fill[color={rgb,255:red,153;green,153;blue,153}] (6.117,6.472) rectangle ++(0.050,0.050);
  \fill[color={rgb,255:red,153;green,153;blue,153}] (5.992,6.477) rectangle ++(0.050,0.050);
  \fill[color={rgb,255:red,153;green,153;blue,153}] (5.866,6.474) rectangle ++(0.050,0.050);
  \fill[color={rgb,255:red,153;green,153;blue,153}] (5.747,6.478) rectangle ++(0.050,0.050);
  \fill[color={rgb,255:red,153;green,153;blue,153}] (5.609,6.451) rectangle ++(0.050,0.050);
  \fill[color={rgb,255:red,153;green,153;blue,153}] (5.514,6.480) rectangle ++(0.050,0.050);
  \fill[color={rgb,255:red,153;green,153;blue,153}] (5.406,6.489) rectangle ++(0.050,0.050);
  \fill[color={rgb,255:red,153;green,153;blue,153}] (5.277,6.460) rectangle ++(0.050,0.050);
  \fill[color={rgb,255:red,153;green,153;blue,153}] (5.181,6.478) rectangle ++(0.050,0.050);
  \fill[color={rgb,255:red,153;green,153;blue,153}] (5.061,6.455) rectangle ++(0.050,0.050);
  \fill[color={rgb,255:red,153;green,153;blue,153}] (4.986,6.504) rectangle ++(0.050,0.050);
  \fill[color={rgb,255:red,153;green,153;blue,153}] (4.870,6.480) rectangle ++(0.050,0.050);
  \fill[color={rgb,255:red,153;green,153;blue,153}] (4.787,6.512) rectangle ++(0.050,0.050);
  \fill[color={rgb,255:red,153;green,153;blue,153}] (4.675,6.486) rectangle ++(0.050,0.050);
  \fill[color={rgb,255:red,153;green,153;blue,153}] (4.578,6.485) rectangle ++(0.050,0.050);
  \fill[color={rgb,255:red,153;green,153;blue,153}] (4.489,6.499) rectangle ++(0.050,0.050);
  \fill[color={rgb,255:red,153;green,153;blue,153}] (4.396,6.501) rectangle ++(0.050,0.050);
  \fill[color={rgb,255:red,153;green,153;blue,153}] (4.303,6.498) rectangle ++(0.050,0.050);
  \fill[color={rgb,255:red,153;green,153;blue,153}] (4.202,6.473) rectangle ++(0.050,0.050);
  \fill[color={rgb,255:red,153;green,153;blue,153}] (4.112,6.472) rectangle ++(0.050,0.050);
  \fill[color={rgb,255:red,153;green,153;blue,153}] (4.030,6.492) rectangle ++(0.050,0.050);
  \fill[color={rgb,255:red,153;green,153;blue,153}] (3.930,6.454) rectangle ++(0.050,0.050);
  \fill[color={rgb,255:red,153;green,153;blue,153}] (3.843,6.450) rectangle ++(0.050,0.050);
  \fill[color={rgb,255:red,153;green,153;blue,153}] (3.768,6.486) rectangle ++(0.050,0.050);
  \fill[color={rgb,255:red,153;green,153;blue,153}] (3.681,6.477) rectangle ++(0.050,0.050);
  \fill[color={rgb,255:red,153;green,153;blue,153}] (3.599,6.484) rectangle ++(0.050,0.050);
  \fill[color={rgb,255:red,153;green,153;blue,153}] (3.515,6.481) rectangle ++(0.050,0.050);
  \fill[color={rgb,255:red,153;green,153;blue,153}] (3.426,6.450) rectangle ++(0.050,0.050);
  \fill[color={rgb,255:red,153;green,153;blue,153}] (3.351,6.481) rectangle ++(0.050,0.050);
  \fill[color={rgb,255:red,153;green,153;blue,153}] (3.270,6.482) rectangle ++(0.050,0.050);
  \fill[color={rgb,255:red,153;green,153;blue,153}] (3.189,6.483) rectangle ++(0.050,0.050);
  \fill[color={rgb,255:red,153;green,153;blue,153}] (3.106,6.462) rectangle ++(0.050,0.050);
  \fill[color={rgb,255:red,153;green,153;blue,153}] (3.027,6.474) rectangle ++(0.050,0.050);
  \fill[color={rgb,255:red,153;green,153;blue,153}] (2.946,6.455) rectangle ++(0.050,0.050);
  \fill[color={rgb,255:red,153;green,153;blue,153}] (2.869,6.478) rectangle ++(0.050,0.050);
  \fill[color={rgb,255:red,153;green,153;blue,153}] (2.790,6.481) rectangle ++(0.050,0.050);
  \fill[color={rgb,255:red,153;green,153;blue,153}] (2.710,6.465) rectangle ++(0.050,0.050);
  \fill[color={rgb,255:red,153;green,153;blue,153}] (2.631,6.455) rectangle ++(0.050,0.050);
  \fill[color={rgb,255:red,153;green,153;blue,153}] (2.554,6.485) rectangle ++(0.050,0.050);
  \fill[color={rgb,255:red,153;green,153;blue,153}] (2.475,6.451) rectangle ++(0.050,0.050);
  \fill[color={rgb,255:red,153;green,153;blue,153}] (2.396,6.482) rectangle ++(0.050,0.050);
  \fill[color={rgb,255:red,153;green,153;blue,153}] (2.317,6.492) rectangle ++(0.050,0.050);
  \fill[color={rgb,255:red,153;green,153;blue,153}] (2.240,6.462) rectangle ++(0.050,0.050);
  \fill[color={rgb,255:red,153;green,153;blue,153}] (2.159,6.498) rectangle ++(0.050,0.050);
  \fill[color={rgb,255:red,153;green,153;blue,153}] (2.081,6.475) rectangle ++(0.050,0.050);
  \fill[color={rgb,255:red,153;green,153;blue,153}] (2.002,6.480) rectangle ++(0.050,0.050);
  \fill[color={rgb,255:red,153;green,153;blue,153}] (1.928,6.431) rectangle ++(0.050,0.050);
  \fill[color={rgb,255:red,153;green,153;blue,153}] (1.841,6.489) rectangle ++(0.050,0.050);
  \fill[color={rgb,255:red,153;green,153;blue,153}] (1.763,6.472) rectangle ++(0.050,0.050);
  \fill[color={rgb,255:red,153;green,153;blue,153}] (1.683,6.469) rectangle ++(0.050,0.050);
  \fill[color={rgb,255:red,153;green,153;blue,153}] (1.600,6.479) rectangle ++(0.050,0.050);
  \fill[color={rgb,255:red,153;green,153;blue,153}] (1.519,6.473) rectangle ++(0.050,0.050);
  \fill[color={rgb,255:red,153;green,153;blue,153}] (1.435,6.479) rectangle ++(0.050,0.050);
  \fill[color={rgb,255:red,153;green,153;blue,153}] (1.356,6.462) rectangle ++(0.050,0.050);
  \fill[color={rgb,255:red,153;green,153;blue,153}] (1.273,6.460) rectangle ++(0.050,0.050);
  \fill[color={rgb,255:red,153;green,153;blue,153}] (1.182,6.485) rectangle ++(0.050,0.050);
  \fill[color={rgb,255:red,153;green,153;blue,153}] (1.104,6.459) rectangle ++(0.050,0.050);
  \fill[color={rgb,255:red,153;green,153;blue,153}] (1.016,6.466) rectangle ++(0.050,0.050);
  \fill[color={rgb,255:red,153;green,153;blue,153}] (0.923,6.481) rectangle ++(0.050,0.050);
  \fill[color={rgb,255:red,153;green,153;blue,153}] (0.836,6.478) rectangle ++(0.050,0.050);
  \fill[color={rgb,255:red,153;green,153;blue,153}] (0.747,6.476) rectangle ++(0.050,0.050);
  \fill[color={rgb,255:red,153;green,153;blue,153}] (0.649,6.494) rectangle ++(0.050,0.050);
  \fill[color={rgb,255:red,153;green,153;blue,153}] (0.565,6.475) rectangle ++(0.050,0.050);
  \fill[color={rgb,255:red,153;green,153;blue,153}] (0.475,6.466) rectangle ++(0.050,0.050);
  \fill[color={rgb,255:red,153;green,153;blue,153}] (0.470,6.276) rectangle ++(0.050,0.050);
  \fill[color={rgb,255:red,153;green,153;blue,153}] (0.477,6.071) rectangle ++(0.050,0.050);
  \fill[color={rgb,255:red,153;green,153;blue,153}] (0.483,5.885) rectangle ++(0.050,0.050);
  \fill[color={rgb,255:red,153;green,153;blue,153}] (0.476,5.735) rectangle ++(0.050,0.050);
  \fill[color={rgb,255:red,153;green,153;blue,153}] (0.472,5.589) rectangle ++(0.050,0.050);
  \fill[color={rgb,255:red,153;green,153;blue,153}] (0.480,5.430) rectangle ++(0.050,0.050);
  \fill[color={rgb,255:red,153;green,153;blue,153}] (0.483,5.290) rectangle ++(0.050,0.050);
  \fill[color={rgb,255:red,153;green,153;blue,153}] (0.471,5.182) rectangle ++(0.050,0.050);
  \fill[color={rgb,255:red,153;green,153;blue,153}] (0.474,5.056) rectangle ++(0.050,0.050);
  \fill[color={rgb,255:red,153;green,153;blue,153}] (0.467,4.952) rectangle ++(0.050,0.050);
  \fill[color={rgb,255:red,153;green,153;blue,153}] (0.476,4.829) rectangle ++(0.050,0.050);
  \fill[color={rgb,255:red,153;green,153;blue,153}] (0.457,4.752) rectangle ++(0.050,0.050);
  \fill[color={rgb,255:red,153;green,153;blue,153}] (0.485,4.616) rectangle ++(0.050,0.050);
  \fill[color={rgb,255:red,153;green,153;blue,153}] (0.487,4.520) rectangle ++(0.050,0.050);
  \fill[color={rgb,255:red,153;green,153;blue,153}] (0.492,4.423) rectangle ++(0.050,0.050);
  \fill[color={rgb,255:red,153;green,153;blue,153}] (0.464,4.372) rectangle ++(0.050,0.050);
  \fill[color={rgb,255:red,153;green,153;blue,153}] (0.463,4.290) rectangle ++(0.050,0.050);
  \fill[color={rgb,255:red,153;green,153;blue,153}] (0.479,4.192) rectangle ++(0.050,0.050);
  \fill[color={rgb,255:red,153;green,153;blue,153}] (0.459,4.137) rectangle ++(0.050,0.050);
  \fill[color={rgb,255:red,153;green,153;blue,153}] (0.489,4.032) rectangle ++(0.050,0.050);
  \fill[color={rgb,255:red,153;green,153;blue,153}] (0.472,3.978) rectangle ++(0.050,0.050);
  \fill[color={rgb,255:red,153;green,153;blue,153}] (0.472,3.909) rectangle ++(0.050,0.050);
  \fill[color={rgb,255:red,153;green,153;blue,153}] (0.473,3.842) rectangle ++(0.050,0.050);
  \fill[color={rgb,255:red,153;green,153;blue,153}] (0.470,3.781) rectangle ++(0.050,0.050);
  \fill[color={rgb,255:red,153;green,153;blue,153}] (0.482,3.707) rectangle ++(0.050,0.050);
  \fill[color={rgb,255:red,153;green,153;blue,153}] (0.468,3.659) rectangle ++(0.050,0.050);
  \fill[color={rgb,255:red,153;green,153;blue,153}] (0.467,3.601) rectangle ++(0.050,0.050);
  \fill[color={rgb,255:red,153;green,153;blue,153}] (0.492,3.524) rectangle ++(0.050,0.050);
  \fill[color={rgb,255:red,153;green,153;blue,153}] (0.463,3.491) rectangle ++(0.050,0.050);
  \fill[color={rgb,255:red,153;green,153;blue,153}] (0.489,3.418) rectangle ++(0.050,0.050);
  \fill[color={rgb,255:red,153;green,153;blue,153}] (0.476,3.375) rectangle ++(0.050,0.050);
  \fill[color={rgb,255:red,153;green,153;blue,153}] (0.457,3.336) rectangle ++(0.050,0.050);
  \fill[color={rgb,255:red,153;green,153;blue,153}] (0.492,3.263) rectangle ++(0.050,0.050);
  \fill[color={rgb,255:red,153;green,153;blue,153}] (0.464,3.231) rectangle ++(0.050,0.050);
  \fill[color={rgb,255:red,153;green,153;blue,153}] (0.466,3.182) rectangle ++(0.050,0.050);
  \fill[color={rgb,255:red,153;green,153;blue,153}] (0.466,3.135) rectangle ++(0.050,0.050);
  \fill[color={rgb,255:red,153;green,153;blue,153}] (0.493,3.074) rectangle ++(0.050,0.050);
  \fill[color={rgb,255:red,153;green,153;blue,153}] (0.473,3.040) rectangle ++(0.050,0.050);
  \fill[color={rgb,255:red,153;green,153;blue,153}] (0.453,3.005) rectangle ++(0.050,0.050);
  \fill[color={rgb,255:red,153;green,153;blue,153}] (0.489,2.944) rectangle ++(0.050,0.050);
  \fill[color={rgb,255:red,153;green,153;blue,153}] (0.472,2.909) rectangle ++(0.050,0.050);
  \fill[color={rgb,255:red,153;green,153;blue,153}] (0.468,2.869) rectangle ++(0.050,0.050);
  \fill[color={rgb,255:red,153;green,153;blue,153}] (0.487,2.819) rectangle ++(0.050,0.050);
  \fill[color={rgb,255:red,153;green,153;blue,153}] (0.497,2.774) rectangle ++(0.050,0.050);
  \fill[color={rgb,255:red,153;green,153;blue,153}] (0.491,2.736) rectangle ++(0.050,0.050);
  \fill[color={rgb,255:red,153;green,153;blue,153}] (0.456,2.710) rectangle ++(0.050,0.050);
  \fill[color={rgb,255:red,153;green,153;blue,153}] (0.462,2.668) rectangle ++(0.050,0.050);
  \fill[color={rgb,255:red,153;green,153;blue,153}] (0.469,2.627) rectangle ++(0.050,0.050);
  \fill[color={rgb,255:red,153;green,153;blue,153}] (0.477,2.586) rectangle ++(0.050,0.050);
  \fill[color={rgb,255:red,153;green,153;blue,153}] (0.478,2.548) rectangle ++(0.050,0.050);
  \fill[color={rgb,255:red,153;green,153;blue,153}] (0.480,2.510) rectangle ++(0.050,0.050);
  \fill[color={rgb,255:red,153;green,153;blue,153}] (0.489,2.470) rectangle ++(0.050,0.050);
  \fill[color={rgb,255:red,153;green,153;blue,153}] (0.474,2.437) rectangle ++(0.050,0.050);
  \fill[color={rgb,255:red,153;green,153;blue,153}] (0.465,2.402) rectangle ++(0.050,0.050);
  \fill[color={rgb,255:red,153;green,153;blue,153}] (0.489,2.361) rectangle ++(0.050,0.050);
  \fill[color={rgb,255:red,153;green,153;blue,153}] (0.474,2.328) rectangle ++(0.050,0.050);
  \fill[color={rgb,255:red,153;green,153;blue,153}] (0.481,2.291) rectangle ++(0.050,0.050);
  \fill[color={rgb,255:red,153;green,153;blue,153}] (0.477,2.256) rectangle ++(0.050,0.050);
  \fill[color={rgb,255:red,153;green,153;blue,153}] (0.464,2.222) rectangle ++(0.050,0.050);
  \fill[color={rgb,255:red,153;green,153;blue,153}] (0.467,2.186) rectangle ++(0.050,0.050);
  \fill[color={rgb,255:red,153;green,153;blue,153}] (0.472,2.150) rectangle ++(0.050,0.050);
  \fill[color={rgb,255:red,153;green,153;blue,153}] (0.503,2.113) rectangle ++(0.050,0.050);
  \fill[color={rgb,255:red,153;green,153;blue,153}] (0.444,2.081) rectangle ++(0.050,0.050);
  \fill[color={rgb,255:red,153;green,153;blue,153}] (0.470,2.045) rectangle ++(0.050,0.050);
  \fill[color={rgb,255:red,153;green,153;blue,153}] (0.483,2.010) rectangle ++(0.050,0.050);
  \fill[color={rgb,255:red,153;green,153;blue,153}] (0.459,1.975) rectangle ++(0.050,0.050);
  \fill[color={rgb,255:red,153;green,153;blue,153}] (0.472,1.940) rectangle ++(0.050,0.050);
  \fill[color={rgb,255:red,153;green,153;blue,153}] (0.478,1.905) rectangle ++(0.050,0.050);
  \fill[color={rgb,255:red,153;green,153;blue,153}] (0.460,1.869) rectangle ++(0.050,0.050);
  \fill[color={rgb,255:red,153;green,153;blue,153}] (0.480,1.835) rectangle ++(0.050,0.050);
  \fill[color={rgb,255:red,153;green,153;blue,153}] (0.466,1.799) rectangle ++(0.050,0.050);
  \fill[color={rgb,255:red,153;green,153;blue,153}] (0.493,1.767) rectangle ++(0.050,0.050);
  \fill[color={rgb,255:red,153;green,153;blue,153}] (0.474,1.729) rectangle ++(0.050,0.050);
  \fill[color={rgb,255:red,153;green,153;blue,153}] (0.509,1.699) rectangle ++(0.050,0.050);
  \fill[color={rgb,255:red,153;green,153;blue,153}] (0.480,1.659) rectangle ++(0.050,0.050);
  \fill[color={rgb,255:red,153;green,153;blue,153}] (0.492,1.625) rectangle ++(0.050,0.050);
  \fill[color={rgb,255:red,153;green,153;blue,153}] (0.483,1.588) rectangle ++(0.050,0.050);
  \fill[color={rgb,255:red,153;green,153;blue,153}] (0.479,1.551) rectangle ++(0.050,0.050);
  \fill[color={rgb,255:red,153;green,153;blue,153}] (0.485,1.516) rectangle ++(0.050,0.050);
  \fill[color={rgb,255:red,153;green,153;blue,153}] (0.442,1.468) rectangle ++(0.050,0.050);
  \fill[color={rgb,255:red,153;green,153;blue,153}] (0.485,1.442) rectangle ++(0.050,0.050);
  \fill[color={rgb,255:red,153;green,153;blue,153}] (0.484,1.404) rectangle ++(0.050,0.050);
  \fill[color={rgb,255:red,153;green,153;blue,153}] (0.475,1.364) rectangle ++(0.050,0.050);
  \fill[color={rgb,255:red,153;green,153;blue,153}] (0.455,1.319) rectangle ++(0.050,0.050);
  \fill[color={rgb,255:red,153;green,153;blue,153}] (0.462,1.282) rectangle ++(0.050,0.050);
  \fill[color={rgb,255:red,153;green,153;blue,153}] (0.459,1.241) rectangle ++(0.050,0.050);
  \fill[color={rgb,255:red,153;green,153;blue,153}] (0.488,1.212) rectangle ++(0.050,0.050);
  \fill[color={rgb,255:red,153;green,153;blue,153}] (0.495,1.175) rectangle ++(0.050,0.050);
  \fill[color={rgb,255:red,153;green,153;blue,153}] (0.491,1.133) rectangle ++(0.050,0.050);
  \fill[color={rgb,255:red,153;green,153;blue,153}] (0.474,1.084) rectangle ++(0.050,0.050);
  \fill[color={rgb,255:red,153;green,153;blue,153}] (0.488,1.048) rectangle ++(0.050,0.050);
  \fill[color={rgb,255:red,153;green,153;blue,153}] (0.485,1.004) rectangle ++(0.050,0.050);
  \fill[color={rgb,255:red,153;green,153;blue,153}] (0.482,0.959) rectangle ++(0.050,0.050);
  \fill[color={rgb,255:red,153;green,153;blue,153}] (0.465,0.906) rectangle ++(0.050,0.050);
  \fill[color={rgb,255:red,153;green,153;blue,153}] (0.474,0.866) rectangle ++(0.050,0.050);
  \fill[color={rgb,255:red,153;green,153;blue,153}] (0.474,0.820) rectangle ++(0.050,0.050);
  \fill[color={rgb,255:red,153;green,153;blue,153}] (0.475,0.773) rectangle ++(0.050,0.050);
  \fill[color={rgb,255:red,153;green,153;blue,153}] (0.483,0.730) rectangle ++(0.050,0.050);
  \fill[color={rgb,255:red,153;green,153;blue,153}] (0.463,0.668) rectangle ++(0.050,0.050);
  \fill[color={rgb,255:red,153;green,153;blue,153}] (0.470,0.622) rectangle ++(0.050,0.050);
  \fill[color={rgb,255:red,153;green,153;blue,153}] (0.486,0.582) rectangle ++(0.050,0.050);
  \fill[color={rgb,255:red,153;green,153;blue,153}] (0.471,0.519) rectangle ++(0.050,0.050);
  \fill[color={rgb,255:red,153;green,153;blue,153}] (0.488,0.478) rectangle ++(0.050,0.050);
  \fill[color={rgb,255:red,153;green,153;blue,153}] (0.543,0.465) rectangle ++(0.050,0.050);
  \fill[color={rgb,255:red,153;green,153;blue,153}] (0.605,0.461) rectangle ++(0.050,0.050);
  \fill[color={rgb,255:red,153;green,153;blue,153}] (0.687,0.475) rectangle ++(0.050,0.050);
  \fill[color={rgb,255:red,153;green,153;blue,153}] (0.740,0.467) rectangle ++(0.050,0.050);
  \fill[color={rgb,255:red,153;green,153;blue,153}] (0.797,0.464) rectangle ++(0.050,0.050);
  \fill[color={rgb,255:red,153;green,153;blue,153}] (0.875,0.483) rectangle ++(0.050,0.050);
  \fill[color={rgb,255:red,153;green,153;blue,153}] (0.904,0.458) rectangle ++(0.050,0.050);
  \fill[color={rgb,255:red,153;green,153;blue,153}] (0.978,0.478) rectangle ++(0.050,0.050);
  \fill[color={rgb,255:red,153;green,153;blue,153}] (1.011,0.459) rectangle ++(0.050,0.050);
  \fill[color={rgb,255:red,153;green,153;blue,153}] (1.074,0.473) rectangle ++(0.050,0.050);
  \fill[color={rgb,255:red,153;green,153;blue,153}] (1.111,0.460) rectangle ++(0.050,0.050);
  \fill[color={rgb,255:red,153;green,153;blue,153}] (1.175,0.480) rectangle ++(0.050,0.050);
  \fill[color={rgb,255:red,153;green,153;blue,153}] (1.223,0.483) rectangle ++(0.050,0.050);
  \fill[color={rgb,255:red,153;green,153;blue,153}] (1.269,0.486) rectangle ++(0.050,0.050);
  \fill[color={rgb,255:red,153;green,153;blue,153}] (1.300,0.470) rectangle ++(0.050,0.050);
  \fill[color={rgb,255:red,153;green,153;blue,153}] (1.341,0.471) rectangle ++(0.050,0.050);
  \fill[color={rgb,255:red,153;green,153;blue,153}] (1.379,0.466) rectangle ++(0.050,0.050);
  \fill[color={rgb,255:red,153;green,153;blue,153}] (1.429,0.481) rectangle ++(0.050,0.050);
  \fill[color={rgb,255:red,153;green,153;blue,153}] (1.467,0.480) rectangle ++(0.050,0.050);
  \fill[color={rgb,255:red,153;green,153;blue,153}] (1.502,0.477) rectangle ++(0.050,0.050);
  \fill[color={rgb,255:red,153;green,153;blue,153}] (1.530,0.463) rectangle ++(0.050,0.050);
  \fill[color={rgb,255:red,153;green,153;blue,153}] (1.576,0.478) rectangle ++(0.050,0.050);
  \fill[color={rgb,255:red,153;green,153;blue,153}] (1.606,0.470) rectangle ++(0.050,0.050);
  \fill[color={rgb,255:red,153;green,153;blue,153}] (1.639,0.466) rectangle ++(0.050,0.050);
  \fill[color={rgb,255:red,153;green,153;blue,153}] (1.678,0.476) rectangle ++(0.050,0.050);
  \fill[color={rgb,255:red,153;green,153;blue,153}] (1.710,0.474) rectangle ++(0.050,0.050);
  \fill[color={rgb,255:red,153;green,153;blue,153}] (1.746,0.479) rectangle ++(0.050,0.050);
  \fill[color={rgb,255:red,153;green,153;blue,153}] (1.772,0.468) rectangle ++(0.050,0.050);
  \fill[color={rgb,255:red,153;green,153;blue,153}] (1.812,0.485) rectangle ++(0.050,0.050);
  \fill[color={rgb,255:red,153;green,153;blue,153}] (1.832,0.460) rectangle ++(0.050,0.050);
  \fill[color={rgb,255:red,153;green,153;blue,153}] (1.856,0.442) rectangle ++(0.050,0.050);
  \fill[color={rgb,255:red,153;green,153;blue,153}] (1.891,0.453) rectangle ++(0.050,0.050);
  \fill[color={rgb,255:red,153;green,153;blue,153}] (1.930,0.477) rectangle ++(0.050,0.050);
  \fill[color={rgb,255:red,153;green,153;blue,153}] (1.961,0.482) rectangle ++(0.050,0.050);
  \fill[color={rgb,255:red,153;green,153;blue,153}] (1.989,0.480) rectangle ++(0.050,0.050);
  \fill[color={rgb,255:red,153;green,153;blue,153}] (2.016,0.473) rectangle ++(0.050,0.050);
  \fill[color={rgb,255:red,153;green,153;blue,153}] (2.042,0.463) rectangle ++(0.050,0.050);
  \fill[color={rgb,255:red,153;green,153;blue,153}] (2.073,0.476) rectangle ++(0.050,0.050);
  \fill[color={rgb,255:red,153;green,153;blue,153}] (2.097,0.457) rectangle ++(0.050,0.050);
  \fill[color={rgb,255:red,153;green,153;blue,153}] (2.132,0.491) rectangle ++(0.050,0.050);
  \fill[color={rgb,255:red,153;green,153;blue,153}] (2.156,0.475) rectangle ++(0.050,0.050);
  \fill[color={rgb,255:red,153;green,153;blue,153}] (2.180,0.459) rectangle ++(0.050,0.050);
  \fill[color={rgb,255:red,153;green,153;blue,153}] (2.209,0.465) rectangle ++(0.050,0.050);
  \fill[color={rgb,255:red,153;green,153;blue,153}] (2.238,0.478) rectangle ++(0.050,0.050);
  \fill[color={rgb,255:red,153;green,153;blue,153}] (2.264,0.471) rectangle ++(0.050,0.050);
  \fill[color={rgb,255:red,153;green,153;blue,153}] (2.291,0.478) rectangle ++(0.050,0.050);
  \fill[color={rgb,255:red,153;green,153;blue,153}] (2.319,0.489) rectangle ++(0.050,0.050);
  \fill[color={rgb,255:red,153;green,153;blue,153}] (2.343,0.461) rectangle ++(0.050,0.050);
  \fill[color={rgb,255:red,153;green,153;blue,153}] (2.370,0.478) rectangle ++(0.050,0.050);
  \fill[color={rgb,255:red,153;green,153;blue,153}] (2.397,0.490) rectangle ++(0.050,0.050);
  \fill[color={rgb,255:red,153;green,153;blue,153}] (2.423,0.497) rectangle ++(0.050,0.050);
  \fill[color={rgb,255:red,153;green,153;blue,153}] (2.449,0.491) rectangle ++(0.050,0.050);
  \fill[color={rgb,255:red,153;green,153;blue,153}] (2.475,0.473) rectangle ++(0.050,0.050);
  \fill[color={rgb,255:red,153;green,153;blue,153}] (2.501,0.496) rectangle ++(0.050,0.050);
  \fill[color={rgb,255:red,153;green,153;blue,153}] (2.528,0.458) rectangle ++(0.050,0.050);
  \fill[color={rgb,255:red,153;green,153;blue,153}] (2.554,0.470) rectangle ++(0.050,0.050);
  \fill[color={rgb,255:red,153;green,153;blue,153}] (2.580,0.480) rectangle ++(0.050,0.050);
  \fill[color={rgb,255:red,153;green,153;blue,153}] (2.604,0.496) rectangle ++(0.050,0.050);
  \fill[color={rgb,255:red,153;green,153;blue,153}] (2.631,0.494) rectangle ++(0.050,0.050);
  \fill[color={rgb,255:red,153;green,153;blue,153}] (2.656,0.499) rectangle ++(0.050,0.050);
  \fill[color={rgb,255:red,153;green,153;blue,153}] (2.683,0.492) rectangle ++(0.050,0.050);
  \fill[color={rgb,255:red,153;green,153;blue,153}] (2.716,0.456) rectangle ++(0.050,0.050);
  \fill[color={rgb,255:red,153;green,153;blue,153}] (2.738,0.486) rectangle ++(0.050,0.050);
  \fill[color={rgb,255:red,153;green,153;blue,153}] (2.772,0.449) rectangle ++(0.050,0.050);
  \fill[color={rgb,255:red,153;green,153;blue,153}] (2.795,0.469) rectangle ++(0.050,0.050);
  \fill[color={rgb,255:red,153;green,153;blue,153}] (2.822,0.473) rectangle ++(0.050,0.050);
  \fill[color={rgb,255:red,153;green,153;blue,153}] (2.853,0.459) rectangle ++(0.050,0.050);
  \fill[color={rgb,255:red,153;green,153;blue,153}] (2.882,0.457) rectangle ++(0.050,0.050);
  \fill[color={rgb,255:red,153;green,153;blue,153}] (2.900,0.492) rectangle ++(0.050,0.050);
  \fill[color={rgb,255:red,153;green,153;blue,153}] (2.934,0.475) rectangle ++(0.050,0.050);
  \fill[color={rgb,255:red,153;green,153;blue,153}] (2.960,0.481) rectangle ++(0.050,0.050);
  \fill[color={rgb,255:red,153;green,153;blue,153}] (2.979,0.511) rectangle ++(0.050,0.050);
  \fill[color={rgb,255:red,153;green,153;blue,153}] (3.038,0.429) rectangle ++(0.050,0.050);
  \fill[color={rgb,255:red,153;green,153;blue,153}] (3.047,0.486) rectangle ++(0.050,0.050);
  \fill[color={rgb,255:red,153;green,153;blue,153}] (3.078,0.482) rectangle ++(0.050,0.050);
  \fill[color={rgb,255:red,153;green,153;blue,153}] (3.108,0.485) rectangle ++(0.050,0.050);
  \fill[color={rgb,255:red,153;green,153;blue,153}] (3.149,0.461) rectangle ++(0.050,0.050);
  \fill[color={rgb,255:red,153;green,153;blue,153}] (3.178,0.467) rectangle ++(0.050,0.050);
  \fill[color={rgb,255:red,153;green,153;blue,153}] (3.204,0.480) rectangle ++(0.050,0.050);
  \fill[color={rgb,255:red,153;green,153;blue,153}] (3.244,0.465) rectangle ++(0.050,0.050);
  \fill[color={rgb,255:red,153;green,153;blue,153}] (3.269,0.483) rectangle ++(0.050,0.050);
  \fill[color={rgb,255:red,153;green,153;blue,153}] (3.301,0.485) rectangle ++(0.050,0.050);
  \fill[color={rgb,255:red,153;green,153;blue,153}] (3.331,0.493) rectangle ++(0.050,0.050);
  \fill[color={rgb,255:red,153;green,153;blue,153}] (3.370,0.486) rectangle ++(0.050,0.050);
  \fill[color={rgb,255:red,153;green,153;blue,153}] (3.419,0.465) rectangle ++(0.050,0.050);
  \fill[color={rgb,255:red,153;green,153;blue,153}] (3.440,0.489) rectangle ++(0.050,0.050);
  \fill[color={rgb,255:red,153;green,153;blue,153}] (3.505,0.448) rectangle ++(0.050,0.050);
  \fill[color={rgb,255:red,153;green,153;blue,153}] (3.523,0.478) rectangle ++(0.050,0.050);
  \fill[color={rgb,255:red,153;green,153;blue,153}] (3.573,0.463) rectangle ++(0.050,0.050);
  \fill[color={rgb,255:red,153;green,153;blue,153}] (3.616,0.460) rectangle ++(0.050,0.050);
  \fill[color={rgb,255:red,153;green,153;blue,153}] (3.643,0.481) rectangle ++(0.050,0.050);
  \fill[color={rgb,255:red,153;green,153;blue,153}] (3.684,0.482) rectangle ++(0.050,0.050);
  \fill[color={rgb,255:red,153;green,153;blue,153}] (3.731,0.478) rectangle ++(0.050,0.050);
  \fill[color={rgb,255:red,153;green,153;blue,153}] (3.763,0.493) rectangle ++(0.050,0.050);
  \fill[color={rgb,255:red,153;green,153;blue,153}] (3.809,0.494) rectangle ++(0.050,0.050);
  \fill[color={rgb,255:red,153;green,153;blue,153}] (3.880,0.468) rectangle ++(0.050,0.050);
  \fill[color={rgb,255:red,153;green,153;blue,153}] (3.912,0.487) rectangle ++(0.050,0.050);
  \fill[color={rgb,255:red,153;green,153;blue,153}] (3.993,0.457) rectangle ++(0.050,0.050);
  \fill[color={rgb,255:red,153;green,153;blue,153}] (4.024,0.479) rectangle ++(0.050,0.050);
  \fill[color={rgb,255:red,153;green,153;blue,153}] (4.057,0.500) rectangle ++(0.050,0.050);
  \fill[color={rgb,255:red,153;green,153;blue,153}] (4.163,0.455) rectangle ++(0.050,0.050);
  \fill[color={rgb,255:red,153;green,153;blue,153}] (4.222,0.457) rectangle ++(0.050,0.050);
  \fill[color={rgb,255:red,153;green,153;blue,153}] (4.243,0.492) rectangle ++(0.050,0.050);
  \fill[color={rgb,255:red,153;green,153;blue,153}] (4.328,0.474) rectangle ++(0.050,0.050);
  \fill[color={rgb,255:red,153;green,153;blue,153}] (4.382,0.485) rectangle ++(0.050,0.050);
  \fill[color={rgb,255:red,153;green,153;blue,153}] (4.454,0.483) rectangle ++(0.050,0.050);
  \fill[color={rgb,255:red,153;green,153;blue,153}] (4.539,0.476) rectangle ++(0.050,0.050);
  \fill[color={rgb,255:red,153;green,153;blue,153}] (4.592,0.492) rectangle ++(0.050,0.050);
  \fill[color={rgb,255:red,153;green,153;blue,153}] (4.709,0.468) rectangle ++(0.050,0.050);
  \fill[color={rgb,255:red,153;green,153;blue,153}] (4.771,0.484) rectangle ++(0.050,0.050);
  \fill[color={rgb,255:red,153;green,153;blue,153}] (4.873,0.476) rectangle ++(0.050,0.050);
  \fill[color={rgb,255:red,153;green,153;blue,153}] (4.970,0.476) rectangle ++(0.050,0.050);
  \fill[color={rgb,255:red,153;green,153;blue,153}] (5.035,0.497) rectangle ++(0.050,0.050);
  \fill[color={rgb,255:red,153;green,153;blue,153}] (5.180,0.476) rectangle ++(0.050,0.050);
  \fill[color={rgb,255:red,153;green,153;blue,153}] (5.291,0.478) rectangle ++(0.050,0.050);
  \fill[color={rgb,255:red,153;green,153;blue,153}] (5.405,0.482) rectangle ++(0.050,0.050);
  \fill[color={rgb,255:red,153;green,153;blue,153}] (5.536,0.482) rectangle ++(0.050,0.050);
  \fill[color={rgb,255:red,153;green,153;blue,153}] (5.696,0.473) rectangle ++(0.050,0.050);
  \fill[color={rgb,255:red,153;green,153;blue,153}] (5.838,0.478) rectangle ++(0.050,0.050);
  \fill[color={rgb,255:red,153;green,153;blue,153}] (6.023,0.469) rectangle ++(0.050,0.050);
  \fill[color={rgb,255:red,153;green,153;blue,153}] (6.194,0.473) rectangle ++(0.050,0.050);
  \fill[color={rgb,255:red,153;green,153;blue,153}] (6.402,0.468) rectangle ++(0.050,0.050);
  \fill[color={rgb,255:red,153;green,153;blue,153}] (6.587,0.479) rectangle ++(0.050,0.050);
  \fill[color={rgb,255:red,153;green,153;blue,153}] (6.837,0.473) rectangle ++(0.050,0.050);
  \fill[color={rgb,255:red,153;green,153;blue,153}] (7.099,0.472) rectangle ++(0.050,0.050);
  \fill[color={rgb,255:red,153;green,153;blue,153}] (7.387,0.473) rectangle ++(0.050,0.050);
  \fill[color={rgb,255:red,153;green,153;blue,153}] (7.727,0.469) rectangle ++(0.050,0.050);
  \fill[color={rgb,255:red,153;green,153;blue,153}] (8.075,0.475) rectangle ++(0.050,0.050);
  \fill[color={rgb,255:red,153;green,153;blue,153}] (8.479,0.478) rectangle ++(0.050,0.050);
  \fill[color={rgb,255:red,153;green,153;blue,153}] (8.970,0.476) rectangle ++(0.050,0.050);
  \fill[color={rgb,255:red,153;green,153;blue,153}] (9.474,0.487) rectangle ++(0.050,0.050);
  \fill[color={rgb,255:red,153;green,153;blue,153}] (9.475,0.614) rectangle ++(0.050,0.050);
  \fill[color={rgb,255:red,153;green,153;blue,153}] (9.494,0.737) rectangle ++(0.050,0.050);
  \fill[color={rgb,255:red,153;green,153;blue,153}] (9.465,0.868) rectangle ++(0.050,0.050);
  \fill[color={rgb,255:red,153;green,153;blue,153}] (9.481,0.990) rectangle ++(0.050,0.050);
  \fill[color={rgb,255:red,153;green,153;blue,153}] (9.486,1.114) rectangle ++(0.050,0.050);
  \fill[color={rgb,255:red,153;green,153;blue,153}] (9.485,1.238) rectangle ++(0.050,0.050);
  \fill[color={rgb,255:red,153;green,153;blue,153}] (9.446,1.365) rectangle ++(0.050,0.050);
  \fill[color={rgb,255:red,153;green,153;blue,153}] (9.457,1.487) rectangle ++(0.050,0.050);
  \fill[color={rgb,255:red,153;green,153;blue,153}] (9.477,1.608) rectangle ++(0.050,0.050);
  \fill[color={rgb,255:red,153;green,153;blue,153}] (9.494,1.730) rectangle ++(0.050,0.050);
  \fill[color={rgb,255:red,153;green,153;blue,153}] (9.484,1.853) rectangle ++(0.050,0.050);
  \fill[color={rgb,255:red,31;green,94;blue,168}] (0.461,0.509) rectangle ++(0.070,0.070);
  \fill[color={rgb,255:red,31;green,94;blue,168}] (0.478,0.468) rectangle ++(0.070,0.070);
  \fill[color={rgb,255:red,31;green,94;blue,168}] (0.533,0.455) rectangle ++(0.070,0.070);
  \fill[color={rgb,255:red,31;green,94;blue,168}] (0.595,0.451) rectangle ++(0.070,0.070);
  \fill[color={rgb,255:red,31;green,94;blue,168}] (0.677,0.465) rectangle ++(0.070,0.070);
  \fill[color={rgb,255:red,31;green,94;blue,168}] (0.730,0.457) rectangle ++(0.070,0.070);
  \fill[color={rgb,255:red,31;green,94;blue,168}] (0.787,0.454) rectangle ++(0.070,0.070);
  \fill[color={rgb,255:red,31;green,94;blue,168}] (0.865,0.473) rectangle ++(0.070,0.070);
  \fill[color={rgb,255:red,31;green,94;blue,168}] (0.894,0.448) rectangle ++(0.070,0.070);
  \fill[color={rgb,255:red,31;green,94;blue,168}] (0.968,0.468) rectangle ++(0.070,0.070);
  \fill[color={rgb,255:red,31;green,94;blue,168}] (1.001,0.449) rectangle ++(0.070,0.070);
  \fill[color={rgb,255:red,31;green,94;blue,168}] (1.064,0.463) rectangle ++(0.070,0.070);
  \fill[color={rgb,255:red,31;green,94;blue,168}] (1.101,0.450) rectangle ++(0.070,0.070);
  \fill[color={rgb,255:red,31;green,94;blue,168}] (1.165,0.470) rectangle ++(0.070,0.070);
  \fill[color={rgb,255:red,31;green,94;blue,168}] (1.213,0.473) rectangle ++(0.070,0.070);
  \fill[color={rgb,255:red,31;green,94;blue,168}] (1.259,0.476) rectangle ++(0.070,0.070);
  \fill[color={rgb,255:red,31;green,94;blue,168}] (1.290,0.460) rectangle ++(0.070,0.070);
  \fill[color={rgb,255:red,31;green,94;blue,168}] (1.331,0.461) rectangle ++(0.070,0.070);
  \fill[color={rgb,255:red,31;green,94;blue,168}] (1.369,0.456) rectangle ++(0.070,0.070);
  \fill[color={rgb,255:red,31;green,94;blue,168}] (1.419,0.471) rectangle ++(0.070,0.070);
  \fill[color={rgb,255:red,31;green,94;blue,168}] (1.457,0.470) rectangle ++(0.070,0.070);
  \fill[color={rgb,255:red,31;green,94;blue,168}] (1.492,0.467) rectangle ++(0.070,0.070);
  \fill[color={rgb,255:red,31;green,94;blue,168}] (1.520,0.453) rectangle ++(0.070,0.070);
  \fill[color={rgb,255:red,31;green,94;blue,168}] (1.566,0.468) rectangle ++(0.070,0.070);
  \fill[color={rgb,255:red,31;green,94;blue,168}] (1.596,0.460) rectangle ++(0.070,0.070);
  \fill[color={rgb,255:red,31;green,94;blue,168}] (1.629,0.456) rectangle ++(0.070,0.070);
  \fill[color={rgb,255:red,31;green,94;blue,168}] (1.668,0.466) rectangle ++(0.070,0.070);
  \fill[color={rgb,255:red,31;green,94;blue,168}] (1.700,0.464) rectangle ++(0.070,0.070);
  \fill[color={rgb,255:red,31;green,94;blue,168}] (1.736,0.469) rectangle ++(0.070,0.070);
  \fill[color={rgb,255:red,31;green,94;blue,168}] (1.762,0.458) rectangle ++(0.070,0.070);
  \fill[color={rgb,255:red,31;green,94;blue,168}] (1.802,0.475) rectangle ++(0.070,0.070);
  \fill[color={rgb,255:red,31;green,94;blue,168}] (1.822,0.450) rectangle ++(0.070,0.070);
  \fill[color={rgb,255:red,31;green,94;blue,168}] (1.846,0.432) rectangle ++(0.070,0.070);
  \fill[color={rgb,255:red,31;green,94;blue,168}] (1.881,0.443) rectangle ++(0.070,0.070);
  \fill[color={rgb,255:red,31;green,94;blue,168}] (1.920,0.467) rectangle ++(0.070,0.070);
  \fill[color={rgb,255:red,31;green,94;blue,168}] (1.951,0.472) rectangle ++(0.070,0.070);
  \fill[color={rgb,255:red,31;green,94;blue,168}] (1.979,0.470) rectangle ++(0.070,0.070);
  \fill[color={rgb,255:red,31;green,94;blue,168}] (2.006,0.463) rectangle ++(0.070,0.070);
  \fill[color={rgb,255:red,31;green,94;blue,168}] (2.032,0.453) rectangle ++(0.070,0.070);
  \fill[color={rgb,255:red,31;green,94;blue,168}] (2.063,0.466) rectangle ++(0.070,0.070);
  \fill[color={rgb,255:red,31;green,94;blue,168}] (2.087,0.447) rectangle ++(0.070,0.070);
  \fill[color={rgb,255:red,31;green,94;blue,168}] (2.122,0.481) rectangle ++(0.070,0.070);
  \fill[color={rgb,255:red,31;green,94;blue,168}] (2.146,0.465) rectangle ++(0.070,0.070);
  \fill[color={rgb,255:red,31;green,94;blue,168}] (2.170,0.449) rectangle ++(0.070,0.070);
  \fill[color={rgb,255:red,31;green,94;blue,168}] (2.199,0.455) rectangle ++(0.070,0.070);
  \fill[color={rgb,255:red,31;green,94;blue,168}] (2.228,0.468) rectangle ++(0.070,0.070);
  \fill[color={rgb,255:red,31;green,94;blue,168}] (2.254,0.461) rectangle ++(0.070,0.070);
  \fill[color={rgb,255:red,31;green,94;blue,168}] (2.281,0.468) rectangle ++(0.070,0.070);
  \fill[color={rgb,255:red,31;green,94;blue,168}] (2.309,0.479) rectangle ++(0.070,0.070);
  \fill[color={rgb,255:red,31;green,94;blue,168}] (2.333,0.451) rectangle ++(0.070,0.070);
  \fill[color={rgb,255:red,31;green,94;blue,168}] (2.360,0.468) rectangle ++(0.070,0.070);
  \fill[color={rgb,255:red,31;green,94;blue,168}] (2.387,0.480) rectangle ++(0.070,0.070);
  \fill[color={rgb,255:red,31;green,94;blue,168}] (2.413,0.487) rectangle ++(0.070,0.070);
  \fill[color={rgb,255:red,31;green,94;blue,168}] (2.439,0.481) rectangle ++(0.070,0.070);
  \fill[color={rgb,255:red,31;green,94;blue,168}] (2.465,0.463) rectangle ++(0.070,0.070);
  \fill[color={rgb,255:red,31;green,94;blue,168}] (2.491,0.486) rectangle ++(0.070,0.070);
  \fill[color={rgb,255:red,31;green,94;blue,168}] (2.518,0.448) rectangle ++(0.070,0.070);
  \fill[color={rgb,255:red,31;green,94;blue,168}] (2.544,0.460) rectangle ++(0.070,0.070);
  \fill[color={rgb,255:red,31;green,94;blue,168}] (2.570,0.470) rectangle ++(0.070,0.070);
  \fill[color={rgb,255:red,31;green,94;blue,168}] (2.594,0.486) rectangle ++(0.070,0.070);
  \fill[color={rgb,255:red,31;green,94;blue,168}] (2.621,0.484) rectangle ++(0.070,0.070);
  \fill[color={rgb,255:red,31;green,94;blue,168}] (2.646,0.489) rectangle ++(0.070,0.070);
  \fill[color={rgb,255:red,31;green,94;blue,168}] (2.673,0.482) rectangle ++(0.070,0.070);
  \fill[color={rgb,255:red,31;green,94;blue,168}] (2.706,0.446) rectangle ++(0.070,0.070);
  \fill[color={rgb,255:red,31;green,94;blue,168}] (2.728,0.476) rectangle ++(0.070,0.070);
  \fill[color={rgb,255:red,31;green,94;blue,168}] (2.762,0.439) rectangle ++(0.070,0.070);
  \fill[color={rgb,255:red,31;green,94;blue,168}] (2.785,0.459) rectangle ++(0.070,0.070);
  \fill[color={rgb,255:red,31;green,94;blue,168}] (2.812,0.463) rectangle ++(0.070,0.070);
  \fill[color={rgb,255:red,31;green,94;blue,168}] (2.843,0.449) rectangle ++(0.070,0.070);
  \fill[color={rgb,255:red,31;green,94;blue,168}] (2.872,0.447) rectangle ++(0.070,0.070);
  \fill[color={rgb,255:red,31;green,94;blue,168}] (2.890,0.482) rectangle ++(0.070,0.070);
  \fill[color={rgb,255:red,31;green,94;blue,168}] (2.924,0.465) rectangle ++(0.070,0.070);
  \fill[color={rgb,255:red,31;green,94;blue,168}] (2.950,0.471) rectangle ++(0.070,0.070);
  \fill[color={rgb,255:red,31;green,94;blue,168}] (2.969,0.501) rectangle ++(0.070,0.070);
  \fill[color={rgb,255:red,31;green,94;blue,168}] (3.028,0.419) rectangle ++(0.070,0.070);
  \fill[color={rgb,255:red,31;green,94;blue,168}] (3.037,0.476) rectangle ++(0.070,0.070);
  \fill[color={rgb,255:red,31;green,94;blue,168}] (3.068,0.472) rectangle ++(0.070,0.070);
  \fill[color={rgb,255:red,31;green,94;blue,168}] (3.098,0.475) rectangle ++(0.070,0.070);
  \fill[color={rgb,255:red,31;green,94;blue,168}] (3.139,0.451) rectangle ++(0.070,0.070);
  \fill[color={rgb,255:red,31;green,94;blue,168}] (3.168,0.457) rectangle ++(0.070,0.070);
  \fill[color={rgb,255:red,31;green,94;blue,168}] (3.194,0.470) rectangle ++(0.070,0.070);
  \fill[color={rgb,255:red,31;green,94;blue,168}] (3.234,0.455) rectangle ++(0.070,0.070);
  \fill[color={rgb,255:red,31;green,94;blue,168}] (3.259,0.473) rectangle ++(0.070,0.070);
  \fill[color={rgb,255:red,31;green,94;blue,168}] (3.291,0.475) rectangle ++(0.070,0.070);
  \fill[color={rgb,255:red,31;green,94;blue,168}] (3.321,0.483) rectangle ++(0.070,0.070);
  \fill[color={rgb,255:red,31;green,94;blue,168}] (3.360,0.476) rectangle ++(0.070,0.070);
  \fill[color={rgb,255:red,31;green,94;blue,168}] (3.409,0.455) rectangle ++(0.070,0.070);
  \fill[color={rgb,255:red,31;green,94;blue,168}] (3.430,0.479) rectangle ++(0.070,0.070);
  \fill[color={rgb,255:red,31;green,94;blue,168}] (3.495,0.438) rectangle ++(0.070,0.070);
  \fill[color={rgb,255:red,31;green,94;blue,168}] (3.513,0.468) rectangle ++(0.070,0.070);
  \fill[color={rgb,255:red,31;green,94;blue,168}] (3.563,0.453) rectangle ++(0.070,0.070);
  \fill[color={rgb,255:red,31;green,94;blue,168}] (3.606,0.450) rectangle ++(0.070,0.070);
  \fill[color={rgb,255:red,31;green,94;blue,168}] (3.633,0.471) rectangle ++(0.070,0.070);
  \fill[color={rgb,255:red,31;green,94;blue,168}] (3.674,0.472) rectangle ++(0.070,0.070);
  \fill[color={rgb,255:red,31;green,94;blue,168}] (3.721,0.468) rectangle ++(0.070,0.070);
  \fill[color={rgb,255:red,31;green,94;blue,168}] (3.753,0.483) rectangle ++(0.070,0.070);
  \fill[color={rgb,255:red,31;green,94;blue,168}] (3.799,0.484) rectangle ++(0.070,0.070);
  \fill[color={rgb,255:red,31;green,94;blue,168}] (3.870,0.458) rectangle ++(0.070,0.070);
  \fill[color={rgb,255:red,31;green,94;blue,168}] (3.902,0.477) rectangle ++(0.070,0.070);
  \fill[color={rgb,255:red,31;green,94;blue,168}] (3.983,0.447) rectangle ++(0.070,0.070);
  \fill[color={rgb,255:red,31;green,94;blue,168}] (4.014,0.469) rectangle ++(0.070,0.070);
  \fill[color={rgb,255:red,31;green,94;blue,168}] (4.047,0.490) rectangle ++(0.070,0.070);
  \fill[color={rgb,255:red,31;green,94;blue,168}] (4.153,0.445) rectangle ++(0.070,0.070);
  \fill[color={rgb,255:red,31;green,94;blue,168}] (4.212,0.447) rectangle ++(0.070,0.070);
  \fill[color={rgb,255:red,31;green,94;blue,168}] (4.233,0.482) rectangle ++(0.070,0.070);
  \fill[color={rgb,255:red,31;green,94;blue,168}] (4.318,0.464) rectangle ++(0.070,0.070);
  \fill[color={rgb,255:red,31;green,94;blue,168}] (4.372,0.475) rectangle ++(0.070,0.070);
  \fill[color={rgb,255:red,31;green,94;blue,168}] (4.444,0.473) rectangle ++(0.070,0.070);
  \fill[color={rgb,255:red,31;green,94;blue,168}] (4.529,0.466) rectangle ++(0.070,0.070);
  \fill[color={rgb,255:red,31;green,94;blue,168}] (4.582,0.482) rectangle ++(0.070,0.070);
  \fill[color={rgb,255:red,31;green,94;blue,168}] (4.699,0.458) rectangle ++(0.070,0.070);
  \fill[color={rgb,255:red,31;green,94;blue,168}] (4.761,0.474) rectangle ++(0.070,0.070);
  \fill[color={rgb,255:red,31;green,94;blue,168}] (4.863,0.466) rectangle ++(0.070,0.070);
  \fill[color={rgb,255:red,31;green,94;blue,168}] (4.960,0.466) rectangle ++(0.070,0.070);
  \fill[color={rgb,255:red,31;green,94;blue,168}] (5.025,0.487) rectangle ++(0.070,0.070);
  \fill[color={rgb,255:red,31;green,94;blue,168}] (5.170,0.466) rectangle ++(0.070,0.070);
  \fill[color={rgb,255:red,31;green,94;blue,168}] (5.281,0.468) rectangle ++(0.070,0.070);
  \fill[color={rgb,255:red,31;green,94;blue,168}] (5.395,0.472) rectangle ++(0.070,0.070);
  \fill[color={rgb,255:red,31;green,94;blue,168}] (5.526,0.472) rectangle ++(0.070,0.070);
  \fill[color={rgb,255:red,31;green,94;blue,168}] (5.686,0.463) rectangle ++(0.070,0.070);
  \fill[color={rgb,255:red,31;green,94;blue,168}] (5.828,0.468) rectangle ++(0.070,0.070);
  \fill[color={rgb,255:red,31;green,94;blue,168}] (6.013,0.459) rectangle ++(0.070,0.070);
  \fill[color={rgb,255:red,31;green,94;blue,168}] (6.184,0.463) rectangle ++(0.070,0.070);
  \fill[color={rgb,255:red,31;green,94;blue,168}] (6.392,0.458) rectangle ++(0.070,0.070);
  \fill[color={rgb,255:red,31;green,94;blue,168}] (6.577,0.469) rectangle ++(0.070,0.070);
  \fill[color={rgb,255:red,31;green,94;blue,168}] (6.827,0.463) rectangle ++(0.070,0.070);
  \fill[color={rgb,255:red,31;green,94;blue,168}] (7.089,0.462) rectangle ++(0.070,0.070);
  \fill[color={rgb,255:red,31;green,94;blue,168}] (7.377,0.463) rectangle ++(0.070,0.070);
  \fill[color={rgb,255:red,31;green,94;blue,168}] (7.717,0.459) rectangle ++(0.070,0.070);
  \fill[color={rgb,255:red,31;green,94;blue,168}] (8.065,0.465) rectangle ++(0.070,0.070);
  \fill[color={rgb,255:red,31;green,94;blue,168}] (8.469,0.468) rectangle ++(0.070,0.070);
  \fill[color={rgb,255:red,31;green,94;blue,168}] (8.960,0.466) rectangle ++(0.070,0.070);
  \fill[color={rgb,255:red,31;green,94;blue,168}] (9.464,0.477) rectangle ++(0.070,0.070);
  \draw[color={rgb,255:red,31;green,94;blue,168}, line width=0.53pt] (0.496,0.499) -- (9.499,0.504);
  \fill[color={rgb,255:red,51;green,161;blue,74}] (0.465,6.456) rectangle ++(0.070,0.070);
  \fill[color={rgb,255:red,51;green,161;blue,74}] (0.460,6.266) rectangle ++(0.070,0.070);
  \fill[color={rgb,255:red,51;green,161;blue,74}] (0.467,6.061) rectangle ++(0.070,0.070);
  \fill[color={rgb,255:red,51;green,161;blue,74}] (0.473,5.875) rectangle ++(0.070,0.070);
  \fill[color={rgb,255:red,51;green,161;blue,74}] (0.466,5.725) rectangle ++(0.070,0.070);
  \fill[color={rgb,255:red,51;green,161;blue,74}] (0.462,5.579) rectangle ++(0.070,0.070);
  \fill[color={rgb,255:red,51;green,161;blue,74}] (0.470,5.420) rectangle ++(0.070,0.070);
  \fill[color={rgb,255:red,51;green,161;blue,74}] (0.473,5.280) rectangle ++(0.070,0.070);
  \fill[color={rgb,255:red,51;green,161;blue,74}] (0.461,5.172) rectangle ++(0.070,0.070);
  \fill[color={rgb,255:red,51;green,161;blue,74}] (0.464,5.046) rectangle ++(0.070,0.070);
  \fill[color={rgb,255:red,51;green,161;blue,74}] (0.457,4.942) rectangle ++(0.070,0.070);
  \fill[color={rgb,255:red,51;green,161;blue,74}] (0.466,4.819) rectangle ++(0.070,0.070);
  \fill[color={rgb,255:red,51;green,161;blue,74}] (0.447,4.742) rectangle ++(0.070,0.070);
  \fill[color={rgb,255:red,51;green,161;blue,74}] (0.475,4.606) rectangle ++(0.070,0.070);
  \fill[color={rgb,255:red,51;green,161;blue,74}] (0.477,4.510) rectangle ++(0.070,0.070);
  \fill[color={rgb,255:red,51;green,161;blue,74}] (0.482,4.413) rectangle ++(0.070,0.070);
  \fill[color={rgb,255:red,51;green,161;blue,74}] (0.454,4.362) rectangle ++(0.070,0.070);
  \fill[color={rgb,255:red,51;green,161;blue,74}] (0.453,4.280) rectangle ++(0.070,0.070);
  \fill[color={rgb,255:red,51;green,161;blue,74}] (0.469,4.182) rectangle ++(0.070,0.070);
  \fill[color={rgb,255:red,51;green,161;blue,74}] (0.449,4.127) rectangle ++(0.070,0.070);
  \fill[color={rgb,255:red,51;green,161;blue,74}] (0.479,4.022) rectangle ++(0.070,0.070);
  \fill[color={rgb,255:red,51;green,161;blue,74}] (0.462,3.968) rectangle ++(0.070,0.070);
  \fill[color={rgb,255:red,51;green,161;blue,74}] (0.462,3.899) rectangle ++(0.070,0.070);
  \fill[color={rgb,255:red,51;green,161;blue,74}] (0.463,3.832) rectangle ++(0.070,0.070);
  \fill[color={rgb,255:red,51;green,161;blue,74}] (0.460,3.771) rectangle ++(0.070,0.070);
  \fill[color={rgb,255:red,51;green,161;blue,74}] (0.472,3.697) rectangle ++(0.070,0.070);
  \fill[color={rgb,255:red,51;green,161;blue,74}] (0.458,3.649) rectangle ++(0.070,0.070);
  \fill[color={rgb,255:red,51;green,161;blue,74}] (0.457,3.591) rectangle ++(0.070,0.070);
  \fill[color={rgb,255:red,51;green,161;blue,74}] (0.482,3.514) rectangle ++(0.070,0.070);
  \fill[color={rgb,255:red,51;green,161;blue,74}] (0.453,3.481) rectangle ++(0.070,0.070);
  \fill[color={rgb,255:red,51;green,161;blue,74}] (0.479,3.408) rectangle ++(0.070,0.070);
  \fill[color={rgb,255:red,51;green,161;blue,74}] (0.466,3.365) rectangle ++(0.070,0.070);
  \fill[color={rgb,255:red,51;green,161;blue,74}] (0.447,3.326) rectangle ++(0.070,0.070);
  \fill[color={rgb,255:red,51;green,161;blue,74}] (0.482,3.253) rectangle ++(0.070,0.070);
  \fill[color={rgb,255:red,51;green,161;blue,74}] (0.454,3.221) rectangle ++(0.070,0.070);
  \fill[color={rgb,255:red,51;green,161;blue,74}] (0.456,3.172) rectangle ++(0.070,0.070);
  \fill[color={rgb,255:red,51;green,161;blue,74}] (0.456,3.125) rectangle ++(0.070,0.070);
  \fill[color={rgb,255:red,51;green,161;blue,74}] (0.483,3.064) rectangle ++(0.070,0.070);
  \fill[color={rgb,255:red,51;green,161;blue,74}] (0.463,3.030) rectangle ++(0.070,0.070);
  \fill[color={rgb,255:red,51;green,161;blue,74}] (0.443,2.995) rectangle ++(0.070,0.070);
  \fill[color={rgb,255:red,51;green,161;blue,74}] (0.479,2.934) rectangle ++(0.070,0.070);
  \fill[color={rgb,255:red,51;green,161;blue,74}] (0.462,2.899) rectangle ++(0.070,0.070);
  \fill[color={rgb,255:red,51;green,161;blue,74}] (0.458,2.859) rectangle ++(0.070,0.070);
  \fill[color={rgb,255:red,51;green,161;blue,74}] (0.477,2.809) rectangle ++(0.070,0.070);
  \fill[color={rgb,255:red,51;green,161;blue,74}] (0.487,2.764) rectangle ++(0.070,0.070);
  \fill[color={rgb,255:red,51;green,161;blue,74}] (0.481,2.726) rectangle ++(0.070,0.070);
  \fill[color={rgb,255:red,51;green,161;blue,74}] (0.446,2.700) rectangle ++(0.070,0.070);
  \fill[color={rgb,255:red,51;green,161;blue,74}] (0.452,2.658) rectangle ++(0.070,0.070);
  \fill[color={rgb,255:red,51;green,161;blue,74}] (0.459,2.617) rectangle ++(0.070,0.070);
  \fill[color={rgb,255:red,51;green,161;blue,74}] (0.467,2.576) rectangle ++(0.070,0.070);
  \fill[color={rgb,255:red,51;green,161;blue,74}] (0.468,2.538) rectangle ++(0.070,0.070);
  \fill[color={rgb,255:red,51;green,161;blue,74}] (0.470,2.500) rectangle ++(0.070,0.070);
  \fill[color={rgb,255:red,51;green,161;blue,74}] (0.479,2.460) rectangle ++(0.070,0.070);
  \fill[color={rgb,255:red,51;green,161;blue,74}] (0.464,2.427) rectangle ++(0.070,0.070);
  \fill[color={rgb,255:red,51;green,161;blue,74}] (0.455,2.392) rectangle ++(0.070,0.070);
  \fill[color={rgb,255:red,51;green,161;blue,74}] (0.479,2.351) rectangle ++(0.070,0.070);
  \fill[color={rgb,255:red,51;green,161;blue,74}] (0.464,2.318) rectangle ++(0.070,0.070);
  \fill[color={rgb,255:red,51;green,161;blue,74}] (0.471,2.281) rectangle ++(0.070,0.070);
  \fill[color={rgb,255:red,51;green,161;blue,74}] (0.467,2.246) rectangle ++(0.070,0.070);
  \fill[color={rgb,255:red,51;green,161;blue,74}] (0.454,2.212) rectangle ++(0.070,0.070);
  \fill[color={rgb,255:red,51;green,161;blue,74}] (0.457,2.176) rectangle ++(0.070,0.070);
  \fill[color={rgb,255:red,51;green,161;blue,74}] (0.462,2.140) rectangle ++(0.070,0.070);
  \fill[color={rgb,255:red,51;green,161;blue,74}] (0.493,2.103) rectangle ++(0.070,0.070);
  \fill[color={rgb,255:red,51;green,161;blue,74}] (0.434,2.071) rectangle ++(0.070,0.070);
  \fill[color={rgb,255:red,51;green,161;blue,74}] (0.460,2.035) rectangle ++(0.070,0.070);
  \fill[color={rgb,255:red,51;green,161;blue,74}] (0.473,2.000) rectangle ++(0.070,0.070);
  \fill[color={rgb,255:red,51;green,161;blue,74}] (0.449,1.965) rectangle ++(0.070,0.070);
  \fill[color={rgb,255:red,51;green,161;blue,74}] (0.462,1.930) rectangle ++(0.070,0.070);
  \fill[color={rgb,255:red,51;green,161;blue,74}] (0.468,1.895) rectangle ++(0.070,0.070);
  \fill[color={rgb,255:red,51;green,161;blue,74}] (0.450,1.859) rectangle ++(0.070,0.070);
  \fill[color={rgb,255:red,51;green,161;blue,74}] (0.470,1.825) rectangle ++(0.070,0.070);
  \fill[color={rgb,255:red,51;green,161;blue,74}] (0.456,1.789) rectangle ++(0.070,0.070);
  \fill[color={rgb,255:red,51;green,161;blue,74}] (0.483,1.757) rectangle ++(0.070,0.070);
  \fill[color={rgb,255:red,51;green,161;blue,74}] (0.464,1.719) rectangle ++(0.070,0.070);
  \fill[color={rgb,255:red,51;green,161;blue,74}] (0.499,1.689) rectangle ++(0.070,0.070);
  \fill[color={rgb,255:red,51;green,161;blue,74}] (0.470,1.649) rectangle ++(0.070,0.070);
  \fill[color={rgb,255:red,51;green,161;blue,74}] (0.482,1.615) rectangle ++(0.070,0.070);
  \fill[color={rgb,255:red,51;green,161;blue,74}] (0.473,1.578) rectangle ++(0.070,0.070);
  \fill[color={rgb,255:red,51;green,161;blue,74}] (0.469,1.541) rectangle ++(0.070,0.070);
  \fill[color={rgb,255:red,51;green,161;blue,74}] (0.475,1.506) rectangle ++(0.070,0.070);
  \fill[color={rgb,255:red,51;green,161;blue,74}] (0.432,1.458) rectangle ++(0.070,0.070);
  \fill[color={rgb,255:red,51;green,161;blue,74}] (0.475,1.432) rectangle ++(0.070,0.070);
  \fill[color={rgb,255:red,51;green,161;blue,74}] (0.474,1.394) rectangle ++(0.070,0.070);
  \fill[color={rgb,255:red,51;green,161;blue,74}] (0.465,1.354) rectangle ++(0.070,0.070);
  \fill[color={rgb,255:red,51;green,161;blue,74}] (0.445,1.309) rectangle ++(0.070,0.070);
  \fill[color={rgb,255:red,51;green,161;blue,74}] (0.452,1.272) rectangle ++(0.070,0.070);
  \fill[color={rgb,255:red,51;green,161;blue,74}] (0.449,1.231) rectangle ++(0.070,0.070);
  \fill[color={rgb,255:red,51;green,161;blue,74}] (0.478,1.202) rectangle ++(0.070,0.070);
  \fill[color={rgb,255:red,51;green,161;blue,74}] (0.485,1.165) rectangle ++(0.070,0.070);
  \fill[color={rgb,255:red,51;green,161;blue,74}] (0.481,1.123) rectangle ++(0.070,0.070);
  \fill[color={rgb,255:red,51;green,161;blue,74}] (0.464,1.074) rectangle ++(0.070,0.070);
  \fill[color={rgb,255:red,51;green,161;blue,74}] (0.478,1.038) rectangle ++(0.070,0.070);
  \fill[color={rgb,255:red,51;green,161;blue,74}] (0.475,0.994) rectangle ++(0.070,0.070);
  \fill[color={rgb,255:red,51;green,161;blue,74}] (0.472,0.949) rectangle ++(0.070,0.070);
  \fill[color={rgb,255:red,51;green,161;blue,74}] (0.455,0.896) rectangle ++(0.070,0.070);
  \fill[color={rgb,255:red,51;green,161;blue,74}] (0.464,0.856) rectangle ++(0.070,0.070);
  \fill[color={rgb,255:red,51;green,161;blue,74}] (0.464,0.810) rectangle ++(0.070,0.070);
  \fill[color={rgb,255:red,51;green,161;blue,74}] (0.465,0.763) rectangle ++(0.070,0.070);
  \fill[color={rgb,255:red,51;green,161;blue,74}] (0.473,0.720) rectangle ++(0.070,0.070);
  \fill[color={rgb,255:red,51;green,161;blue,74}] (0.453,0.658) rectangle ++(0.070,0.070);
  \fill[color={rgb,255:red,51;green,161;blue,74}] (0.460,0.612) rectangle ++(0.070,0.070);
  \fill[color={rgb,255:red,51;green,161;blue,74}] (0.476,0.572) rectangle ++(0.070,0.070);
  \draw[color={rgb,255:red,51;green,161;blue,74}, line width=0.53pt] (0.499,6.491) -- (0.502,0.607);
  \fill[color={rgb,255:red,242;green,140;blue,26}] (7.815,6.455) rectangle ++(0.070,0.070);
  \fill[color={rgb,255:red,242;green,140;blue,26}] (7.652,6.474) rectangle ++(0.070,0.070);
  \fill[color={rgb,255:red,242;green,140;blue,26}] (7.471,6.472) rectangle ++(0.070,0.070);
  \fill[color={rgb,255:red,242;green,140;blue,26}] (7.286,6.461) rectangle ++(0.070,0.070);
  \fill[color={rgb,255:red,242;green,140;blue,26}] (7.103,6.444) rectangle ++(0.070,0.070);
  \fill[color={rgb,255:red,242;green,140;blue,26}] (6.939,6.439) rectangle ++(0.070,0.070);
  \fill[color={rgb,255:red,242;green,140;blue,26}] (6.818,6.472) rectangle ++(0.070,0.070);
  \fill[color={rgb,255:red,242;green,140;blue,26}] (6.678,6.483) rectangle ++(0.070,0.070);
  \fill[color={rgb,255:red,242;green,140;blue,26}] (6.500,6.447) rectangle ++(0.070,0.070);
  \fill[color={rgb,255:red,242;green,140;blue,26}] (6.376,6.464) rectangle ++(0.070,0.070);
  \fill[color={rgb,255:red,242;green,140;blue,26}] (6.245,6.470) rectangle ++(0.070,0.070);
  \fill[color={rgb,255:red,242;green,140;blue,26}] (6.107,6.462) rectangle ++(0.070,0.070);
  \fill[color={rgb,255:red,242;green,140;blue,26}] (5.982,6.467) rectangle ++(0.070,0.070);
  \fill[color={rgb,255:red,242;green,140;blue,26}] (5.856,6.464) rectangle ++(0.070,0.070);
  \fill[color={rgb,255:red,242;green,140;blue,26}] (5.737,6.468) rectangle ++(0.070,0.070);
  \fill[color={rgb,255:red,242;green,140;blue,26}] (5.599,6.441) rectangle ++(0.070,0.070);
  \fill[color={rgb,255:red,242;green,140;blue,26}] (5.504,6.470) rectangle ++(0.070,0.070);
  \fill[color={rgb,255:red,242;green,140;blue,26}] (5.396,6.479) rectangle ++(0.070,0.070);
  \fill[color={rgb,255:red,242;green,140;blue,26}] (5.267,6.450) rectangle ++(0.070,0.070);
  \fill[color={rgb,255:red,242;green,140;blue,26}] (5.171,6.468) rectangle ++(0.070,0.070);
  \fill[color={rgb,255:red,242;green,140;blue,26}] (5.051,6.445) rectangle ++(0.070,0.070);
  \fill[color={rgb,255:red,242;green,140;blue,26}] (4.976,6.494) rectangle ++(0.070,0.070);
  \fill[color={rgb,255:red,242;green,140;blue,26}] (4.860,6.470) rectangle ++(0.070,0.070);
  \fill[color={rgb,255:red,242;green,140;blue,26}] (4.777,6.502) rectangle ++(0.070,0.070);
  \fill[color={rgb,255:red,242;green,140;blue,26}] (4.665,6.476) rectangle ++(0.070,0.070);
  \fill[color={rgb,255:red,242;green,140;blue,26}] (4.568,6.475) rectangle ++(0.070,0.070);
  \fill[color={rgb,255:red,242;green,140;blue,26}] (4.479,6.489) rectangle ++(0.070,0.070);
  \fill[color={rgb,255:red,242;green,140;blue,26}] (4.386,6.491) rectangle ++(0.070,0.070);
  \fill[color={rgb,255:red,242;green,140;blue,26}] (4.293,6.488) rectangle ++(0.070,0.070);
  \fill[color={rgb,255:red,242;green,140;blue,26}] (4.192,6.463) rectangle ++(0.070,0.070);
  \fill[color={rgb,255:red,242;green,140;blue,26}] (4.102,6.462) rectangle ++(0.070,0.070);
  \fill[color={rgb,255:red,242;green,140;blue,26}] (4.020,6.482) rectangle ++(0.070,0.070);
  \fill[color={rgb,255:red,242;green,140;blue,26}] (3.920,6.444) rectangle ++(0.070,0.070);
  \fill[color={rgb,255:red,242;green,140;blue,26}] (3.833,6.440) rectangle ++(0.070,0.070);
  \fill[color={rgb,255:red,242;green,140;blue,26}] (3.758,6.476) rectangle ++(0.070,0.070);
  \fill[color={rgb,255:red,242;green,140;blue,26}] (3.671,6.467) rectangle ++(0.070,0.070);
  \fill[color={rgb,255:red,242;green,140;blue,26}] (3.589,6.474) rectangle ++(0.070,0.070);
  \fill[color={rgb,255:red,242;green,140;blue,26}] (3.505,6.471) rectangle ++(0.070,0.070);
  \fill[color={rgb,255:red,242;green,140;blue,26}] (3.416,6.440) rectangle ++(0.070,0.070);
  \fill[color={rgb,255:red,242;green,140;blue,26}] (3.341,6.471) rectangle ++(0.070,0.070);
  \fill[color={rgb,255:red,242;green,140;blue,26}] (3.260,6.472) rectangle ++(0.070,0.070);
  \fill[color={rgb,255:red,242;green,140;blue,26}] (3.179,6.473) rectangle ++(0.070,0.070);
  \fill[color={rgb,255:red,242;green,140;blue,26}] (3.096,6.452) rectangle ++(0.070,0.070);
  \fill[color={rgb,255:red,242;green,140;blue,26}] (3.017,6.464) rectangle ++(0.070,0.070);
  \fill[color={rgb,255:red,242;green,140;blue,26}] (2.936,6.445) rectangle ++(0.070,0.070);
  \fill[color={rgb,255:red,242;green,140;blue,26}] (2.859,6.468) rectangle ++(0.070,0.070);
  \fill[color={rgb,255:red,242;green,140;blue,26}] (2.780,6.471) rectangle ++(0.070,0.070);
  \fill[color={rgb,255:red,242;green,140;blue,26}] (2.700,6.455) rectangle ++(0.070,0.070);
  \fill[color={rgb,255:red,242;green,140;blue,26}] (2.621,6.445) rectangle ++(0.070,0.070);
  \fill[color={rgb,255:red,242;green,140;blue,26}] (2.544,6.475) rectangle ++(0.070,0.070);
  \fill[color={rgb,255:red,242;green,140;blue,26}] (2.465,6.441) rectangle ++(0.070,0.070);
  \fill[color={rgb,255:red,242;green,140;blue,26}] (2.386,6.472) rectangle ++(0.070,0.070);
  \fill[color={rgb,255:red,242;green,140;blue,26}] (2.307,6.482) rectangle ++(0.070,0.070);
  \fill[color={rgb,255:red,242;green,140;blue,26}] (2.230,6.452) rectangle ++(0.070,0.070);
  \fill[color={rgb,255:red,242;green,140;blue,26}] (2.149,6.488) rectangle ++(0.070,0.070);
  \fill[color={rgb,255:red,242;green,140;blue,26}] (2.071,6.465) rectangle ++(0.070,0.070);
  \fill[color={rgb,255:red,242;green,140;blue,26}] (1.992,6.470) rectangle ++(0.070,0.070);
  \fill[color={rgb,255:red,242;green,140;blue,26}] (1.918,6.421) rectangle ++(0.070,0.070);
  \fill[color={rgb,255:red,242;green,140;blue,26}] (1.831,6.479) rectangle ++(0.070,0.070);
  \fill[color={rgb,255:red,242;green,140;blue,26}] (1.753,6.462) rectangle ++(0.070,0.070);
  \fill[color={rgb,255:red,242;green,140;blue,26}] (1.673,6.459) rectangle ++(0.070,0.070);
  \fill[color={rgb,255:red,242;green,140;blue,26}] (1.590,6.469) rectangle ++(0.070,0.070);
  \fill[color={rgb,255:red,242;green,140;blue,26}] (1.509,6.463) rectangle ++(0.070,0.070);
  \fill[color={rgb,255:red,242;green,140;blue,26}] (1.425,6.469) rectangle ++(0.070,0.070);
  \fill[color={rgb,255:red,242;green,140;blue,26}] (1.346,6.452) rectangle ++(0.070,0.070);
  \fill[color={rgb,255:red,242;green,140;blue,26}] (1.263,6.450) rectangle ++(0.070,0.070);
  \fill[color={rgb,255:red,242;green,140;blue,26}] (1.172,6.475) rectangle ++(0.070,0.070);
  \fill[color={rgb,255:red,242;green,140;blue,26}] (1.094,6.449) rectangle ++(0.070,0.070);
  \fill[color={rgb,255:red,242;green,140;blue,26}] (1.006,6.456) rectangle ++(0.070,0.070);
  \fill[color={rgb,255:red,242;green,140;blue,26}] (0.913,6.471) rectangle ++(0.070,0.070);
  \fill[color={rgb,255:red,242;green,140;blue,26}] (0.826,6.468) rectangle ++(0.070,0.070);
  \fill[color={rgb,255:red,242;green,140;blue,26}] (0.737,6.466) rectangle ++(0.070,0.070);
  \fill[color={rgb,255:red,242;green,140;blue,26}] (0.639,6.484) rectangle ++(0.070,0.070);
  \fill[color={rgb,255:red,242;green,140;blue,26}] (0.555,6.465) rectangle ++(0.070,0.070);
  \draw[color={rgb,255:red,242;green,140;blue,26}, line width=0.53pt] (7.850,6.501) -- (0.590,6.500);
  \fill[color={rgb,255:red,140;green,76;blue,178}] (9.474,1.965) rectangle ++(0.070,0.070);
  \fill[color={rgb,255:red,140;green,76;blue,178}] (9.452,2.087) rectangle ++(0.070,0.070);
  \fill[color={rgb,255:red,140;green,76;blue,178}] (9.455,2.209) rectangle ++(0.070,0.070);
  \fill[color={rgb,255:red,140;green,76;blue,178}] (9.479,2.333) rectangle ++(0.070,0.070);
  \fill[color={rgb,255:red,140;green,76;blue,178}] (9.476,2.455) rectangle ++(0.070,0.070);
  \fill[color={rgb,255:red,140;green,76;blue,178}] (9.472,2.578) rectangle ++(0.070,0.070);
  \fill[color={rgb,255:red,140;green,76;blue,178}] (9.467,2.701) rectangle ++(0.070,0.070);
  \fill[color={rgb,255:red,140;green,76;blue,178}] (9.456,2.823) rectangle ++(0.070,0.070);
  \fill[color={rgb,255:red,140;green,76;blue,178}] (9.505,2.954) rectangle ++(0.070,0.070);
  \fill[color={rgb,255:red,140;green,76;blue,178}] (9.484,3.077) rectangle ++(0.070,0.070);
  \fill[color={rgb,255:red,140;green,76;blue,178}] (9.455,3.198) rectangle ++(0.070,0.070);
  \fill[color={rgb,255:red,140;green,76;blue,178}] (9.472,3.327) rectangle ++(0.070,0.070);
  \fill[color={rgb,255:red,140;green,76;blue,178}] (9.490,3.458) rectangle ++(0.070,0.070);
  \fill[color={rgb,255:red,140;green,76;blue,178}] (9.480,3.585) rectangle ++(0.070,0.070);
  \fill[color={rgb,255:red,140;green,76;blue,178}] (9.471,3.712) rectangle ++(0.070,0.070);
  \fill[color={rgb,255:red,140;green,76;blue,178}] (9.501,3.850) rectangle ++(0.070,0.070);
  \fill[color={rgb,255:red,140;green,76;blue,178}] (9.472,3.974) rectangle ++(0.070,0.070);
  \fill[color={rgb,255:red,140;green,76;blue,178}] (9.459,4.103) rectangle ++(0.070,0.070);
  \fill[color={rgb,255:red,140;green,76;blue,178}] (9.468,4.241) rectangle ++(0.070,0.070);
  \fill[color={rgb,255:red,140;green,76;blue,178}] (9.450,4.370) rectangle ++(0.070,0.070);
  \fill[color={rgb,255:red,140;green,76;blue,178}] (9.465,0.604) rectangle ++(0.070,0.070);
  \fill[color={rgb,255:red,140;green,76;blue,178}] (9.484,0.727) rectangle ++(0.070,0.070);
  \fill[color={rgb,255:red,140;green,76;blue,178}] (9.455,0.858) rectangle ++(0.070,0.070);
  \fill[color={rgb,255:red,140;green,76;blue,178}] (9.471,0.980) rectangle ++(0.070,0.070);
  \fill[color={rgb,255:red,140;green,76;blue,178}] (9.476,1.104) rectangle ++(0.070,0.070);
  \fill[color={rgb,255:red,140;green,76;blue,178}] (9.475,1.228) rectangle ++(0.070,0.070);
  \fill[color={rgb,255:red,140;green,76;blue,178}] (9.436,1.355) rectangle ++(0.070,0.070);
  \fill[color={rgb,255:red,140;green,76;blue,178}] (9.447,1.477) rectangle ++(0.070,0.070);
  \fill[color={rgb,255:red,140;green,76;blue,178}] (9.467,1.598) rectangle ++(0.070,0.070);
  \fill[color={rgb,255:red,140;green,76;blue,178}] (9.484,1.720) rectangle ++(0.070,0.070);
  \fill[color={rgb,255:red,140;green,76;blue,178}] (9.474,1.843) rectangle ++(0.070,0.070);
  \draw[color={rgb,255:red,140;green,76;blue,178}, line width=0.53pt] (9.501,0.639) -- (9.509,4.405);
  \draw[color={rgb,255:red,199;green,15;blue,46}, line width=0.88pt, line join=round] (0.502,0.499) -- (9.501,0.504) -- (9.514,6.501) -- (0.499,6.500) -- (0.502,0.499);
  \fill[color={rgb,255:red,199;green,15;blue,46}] (0.502,0.499) circle (0.075);
  \fill[color={rgb,255:red,199;green,15;blue,46}] (9.501,0.504) circle (0.075);
  \fill[color={rgb,255:red,199;green,15;blue,46}] (9.514,6.501) circle (0.075);
  \fill[color={rgb,255:red,199;green,15;blue,46}] (0.499,6.500) circle (0.075);
  \node[color={rgb,255:red,255;green,255;blue,255}, inner sep=0pt] at (5.001,0.502) {9.00 m};
  \node[color={rgb,255:red,255;green,255;blue,255}, inner sep=0pt] at (9.507,3.503) {6.00 m};
  \node[color={rgb,255:red,255;green,255;blue,255}, inner sep=0pt] at (5.006,6.500) {9.01 m};
  \node[color={rgb,255:red,255;green,255;blue,255}, inner sep=0pt] at (0.500,3.499) {6.00 m};
\end{tikzpicture}
}\\[2pt]
    {\small\thema{-p hesse:=true}, your \thema{fit\_line\_hesse}}
  \end{minipage}
  \par\smallskip{\footnotesize\color{ohmgrey}terminal 2:}
  \lstinputlisting[basicstyle=\ttfamily\scriptsize, aboveskip=2pt, belowskip=0pt]{figures/term_z1.txt}
\end{seebox}

Why does RANSAC still find the vertical walls with \thema{line\_through}, which only gives \(m\) and \(t\)?

\subsection*{Z.2: how exact is the circle?}
The node evaluates a new scan every 2\,s. Note the radius of ten evaluations and compute its mean and
\(\sigma\) with your \thema{mean\_and\_std}. Then compare with the circle seen from closer: how would you
expect \(\sigma\) of the radius to change with the distance to the object, and why?

\end{document}
